
\documentclass[12pt,a4paper]{report}
\usepackage{tocloft}

% Right-align page numbers with dot leaders
\renewcommand{\cftchapleader}{\cftdotfill{\cftdotsep}}
\renewcommand{\cftsecleader}{\cftdotfill{\cftdotsep}}
\renewcommand{\cftsubsecleader}{\cftdotfill{\cftdotsep}}
\usepackage[a4paper,margin=1in]{geometry}
\usepackage{setspace}
\usepackage{graphicx}
\usepackage{float}
\usepackage[nottoc]{tocbibind}

\usepackage{booktabs}
\usepackage{array}
\usepackage{longtable}
\usepackage{amsmath}
\usepackage{booktabs}
\usepackage{tabularx}
\usepackage{array}

\newcolumntype{Y}{>{\centering\arraybackslash}X}
\usepackage{placeins}
\usepackage{amssymb}
\usepackage[backend=biber,style=ieee]{biblatex}
\addbibresource{ref.bib}
\usepackage[hidelinks]{hyperref}
\hypersetup{colorlinks=false,pdfborder={0 0 0}}
\usepackage{caption}
\usepackage{enumitem}
\usepackage{listings}
\usepackage{booktabs}
\usepackage{tabularx}
\usepackage{array}
\usepackage{xcolor}
\usepackage{tocloft}
\usepackage[backend=biber,style=ieee]{biblatex}
\addbibresource{ref.bib}
\newcommand{\listequationsname}{List of Equations}
\newlistof{myequations}{equ}{\listequationsname}
\graphicspath{{images/}{figures/}{outputs/}}
\setcounter{tocdepth}{2}
\setcounter{secnumdepth}{2}
\setcounter{tocdepth}{2}
\onehalfspacing
\setlength{\parindent}{0pt}
\setlength{\parskip}{6pt}

\lstset{
    basicstyle=\ttfamily\small,
    breaklines=true,
    frame=single,
    numbers=left,
    numberstyle=\tiny,
    keywordstyle=\color{blue},
    commentstyle=\color{gray},
    stringstyle=\color{purple}
}

\begin{document}

\begin{titlepage}
    \centering
    \vspace*{0.6cm}
    \includegraphics[width=0.55\textwidth]{images/logo.png}\par
    \vspace{1.0cm}
    {\Huge\bfseries A Hybrid SecureBERT-BiLSTM-Attention Model for Cyber Threat Intelligence Classification\par}
    \vspace{1.0cm}
    {\Large\bfseries Saikiran Achuta\par}
    \vspace{0.3cm}
    {\large Student Number: 3196603\par}
    \vspace{0.6cm}
    {\large Programme: MSc in Computing Science\par}
    \vspace{0.6cm}
    {\large Submitted in partial fulfilment of the requirements for the degree of\par}
    {\large  of Science in Computing science\par}
    {\large Griffith College Dublin\par}
    \vspace{0.6cm}
    {\large Supervisor: Abubakr Siddig\par}
    \vspace{0.4cm}
    {\large September 2026\par}
\end{titlepage}

\pagenumbering{roman}
\chapter*{Disclaimer}
\addcontentsline{toc}{chapter}{Disclaimer}

I hereby declare that this work is original and that I have not submitted it for assessment for an academic purpose anywhere else or for the award of any academic qualification at this or any other academic institution except as a part of the programme of study leading to the Degree of Master of Science in Computing Science at Griffith College Dublin.

\vspace{1cm}
 Saikiran Achuta \hfill 08-09-2026
 
\chapter*{Acknowledgement}
\addcontentsline{toc}{chapter}{Acknowledgement}


I would like to express my gratitude to everyone who had contributed to the completion of this dissertation project “A Hybrid SecureBERT-BiLSTM-Attention Model for Cyber Threat Intelligence Classification”.

I would like to thank my project supervisor for his guidance, feedback and support throughout the research and development process, especially the feedback provided on the Week 4 submission which guided the direction of the final model development, evaluation and robustness testing presented in this report.

I would like to thank the Computing Science faculty and staff at Griffith College for providing me with a space to do applied research and independent problem solving, as well as the opensource communities of Python, pandas, NumPy, Scikitlearn, NLTK, TensorFlow/Keras, joblib, Matplotlib, and Seaborn. The tools enabled the system to be practically developed and validated end-to-end.

Finally, I would like to thank my family and friends for their patience, encouragement and motivation during this academic process.

\cleardoublepage
\tableofcontents
\cleardoublepage

\listoffigures
\cleardoublepage

\listoftables
\cleardoublepage
% \cleardoublepage
% \hyperref[abstract]{Abstract \hyperref[abstract]{10}}

% \hyperref[chapter-1-introduction]{Chapter 1 Introduction
% \hyperref[chapter-1-introduction]{1}}

% \hyperref[_Toc239502922]{\textbf{1.1 Overview}
% \hyperref[_Toc239502922]{1}}

% \hyperref[_Toc239502923]{\textbf{1.2 Background}
% \hyperref[_Toc239502923]{1}}

% \hyperref[_Toc239502924]{\textbf{1.3 Rationale}
% \hyperref[_Toc239502924]{2}}

% \hyperref[_Toc239502925]{\textbf{1.4 Problem Statement}
% \hyperref[_Toc239502925]{2}}

% \hyperref[_Toc239502926]{\textbf{1.5 Research Aim}
% \hyperref[_Toc239502926]{3}}

% \hyperref[_Toc239502927]{\textbf{1.6 Research Objectives}
% \hyperref[_Toc239502927]{3}}

% \hyperref[_Toc239502928]{\textbf{1.7 Research Questions}
% \hyperref[_Toc239502928]{3}}

% \hyperref[_Toc239502929]{\textbf{1.8 Detailed Technical Justification}
% \hyperref[_Toc239502929]{4}}

% \hyperref[_Toc239502930]{\textbf{1.9 Extended Technical Discussion}
% \hyperref[_Toc239502930]{4}}

% \hyperref[_Toc239502931]{\textbf{1.10 Research Significance}
% \hyperref[_Toc239502931]{5}}

% \hyperref[_Toc239502932]{\textbf{1.11 Research Structure}
% \hyperref[_Toc239502932]{5}}

% \hyperref[chapter-2-literature-review]{Chapter 2 Literature Review
% \hyperref[chapter-2-literature-review]{7}}

% \hyperref[introduction]{2.1 Introduction \hyperref[introduction]{7}}

% \hyperref[foundational-nlp-architectures]{2.2 Foundational NLP
% Architectures \hyperref[foundational-nlp-architectures]{8}}

% \hyperref[_Toc239502936]{\textbf{2.3 Domain-Adapted Language Models for
% Cybersecurity} \hyperref[_Toc239502936]{9}}

% \hyperref[_Toc239502937]{\textbf{2.4 Hybrid Transformer--Recurrent
% Architectures for Threat Text Classification}
% \hyperref[_Toc239502937]{10}}

% \hyperref[_Toc239502938]{\textbf{2.5 TTP and Entity/Relation Extraction
% from CTI Text} \hyperref[_Toc239502938]{11}}

% \hyperref[large-language-models-in-cti-and-reliability-concerns]{2.6
% Large Language Models in CTI and Reliability Concerns
% \hyperref[large-language-models-in-cti-and-reliability-concerns]{13}}

% \hyperref[identified-research-gaps]{2.7 Identified Research Gaps
% \hyperref[identified-research-gaps]{13}}

% \hyperref[chapter-3-methodology]{Chapter 3 Methodology
% \hyperref[chapter-3-methodology]{16}}

% \hyperref[overview]{3.1 Overview \hyperref[overview]{16}}

% \hyperref[dataset-and-data-structure]{3.2 Dataset and Data Structure
% \hyperref[dataset-and-data-structure]{17}}

% \hyperref[data-cleaning-and-preprocessing]{3.3 Data Cleaning and
% Preprocessing \hyperref[data-cleaning-and-preprocessing]{17}}

% \hyperref[construction-of-the-classification-target]{3.4 Construction of
% the Classification Target
% \hyperref[construction-of-the-classification-target]{18}}

% \hyperref[entity-aware-exploratory-features]{3.5 Entity-Aware
% Exploratory Features \hyperref[entity-aware-exploratory-features]{19}}

% \hyperref[text-level-statistical-features]{3.6 Text-Level Statistical
% Features \hyperref[text-level-statistical-features]{20}}

% \hyperref[exploratory-data-analysis]{3.7 Exploratory Data Analysis
% \hyperref[exploratory-data-analysis]{20}}

% \hyperref[model-training-and-validation]{3.8 Model Training and
% Validation \hyperref[model-training-and-validation]{23}}

% \hyperref[classical-machine-learning-baselines]{3.9 Classical Machine
% Learning Baselines \hyperref[classical-machine-learning-baselines]{24}}

% \hyperref[bilstm-deep-learning-baseline]{3.10 BiLSTM Deep Learning
% Baseline \hyperref[bilstm-deep-learning-baseline]{24}}

% \hyperref[transformer-and-hybrid-architecture]{3.11 Transformer and
% Hybrid Architecture \hyperref[transformer-and-hybrid-architecture]{25}}

% \hyperref[evaluation-and-error-analysis]{3.12 Evaluation and Error
% Analysis \hyperref[evaluation-and-error-analysis]{25}}

% \hyperref[reproducibility-and-experimental-workflow]{3.13
% Reproducibility and Experimental Workflow
% \hyperref[reproducibility-and-experimental-workflow]{26}}

% \hyperref[methodological-considerations-and-scope]{3.14 Methodological
% Considerations and Scope
% \hyperref[methodological-considerations-and-scope]{26}}

% \hyperref[streamlit-dashboard-integration]{3.15 Streamlit Dashboard
% Integration \hyperref[streamlit-dashboard-integration]{27}}

% \hyperref[summary]{3.16 Summary \hyperref[summary]{27}}

% \hyperref[chapter-4-system-analysis-and-design]{Chapter 4 System
% Analysis and Design \hyperref[chapter-4-system-analysis-and-design]{28}}

% \hyperref[overview-of-the-system-architecture]{4.1 Overview of the
% System Architecture \hyperref[overview-of-the-system-architecture]{28}}

% \hyperref[data-layer]{4.2 Data Layer \hyperref[data-layer]{29}}

% \hyperref[preprocessing-and-label-construction-module]{4.3 Preprocessing
% and Label Construction Module
% \hyperref[preprocessing-and-label-construction-module]{29}}

% \hyperref[feature-engineering-module]{4.4 Feature Engineering Module
% \hyperref[feature-engineering-module]{29}}

% \hyperref[exploratory-analysis-module]{4.5 Exploratory Analysis Module
% \hyperref[exploratory-analysis-module]{30}}

% \hyperref[baseline-model-architectures]{4.6 Baseline Model Architectures
% \hyperref[baseline-model-architectures]{30}}

% \hyperref[proposed-hybrid-architecture]{4.7 Proposed Hybrid Architecture
% \hyperref[proposed-hybrid-architecture]{30}}

% \hyperref[evaluation-harness]{4.8 Evaluation Harness
% \hyperref[evaluation-harness]{31}}

% \hyperref[design-rationale-summary]{4.9 Design Rationale Summary
% \hyperref[design-rationale-summary]{31}}

% \hyperref[summary-1]{4.10 Summary \hyperref[summary-1]{31}}

% \hyperref[chapter-5-implementation]{Chapter 5 Implementation
% \hyperref[chapter-5-implementation]{32}}

% \hyperref[technology-stack]{5.1 Technology Stack
% \hyperref[technology-stack]{32}}

% \hyperref[data-preparation-implementation]{5.2 Data Preparation
% Implementation \hyperref[data-preparation-implementation]{32}}

% \hyperref[label-construction-implementation]{5.3 Label Construction
% Implementation \hyperref[label-construction-implementation]{33}}

% \hyperref[feature-engineering-implementation]{5.4 Feature Engineering
% Implementation \hyperref[feature-engineering-implementation]{33}}

% \hyperref[baseline-model-implementation]{5.5 Baseline Model
% Implementation \hyperref[baseline-model-implementation]{34}}

% \hyperref[hybrid-model-implementation]{5.6 Hybrid Model Implementation
% \hyperref[hybrid-model-implementation]{34}}

% \hyperref[evaluation-and-reporting-implementation]{5.7 Evaluation and
% Reporting Implementation
% \hyperref[evaluation-and-reporting-implementation]{35}}

% \hyperref[reproducibility-measures]{5.8 Reproducibility Measures
% \hyperref[reproducibility-measures]{35}}

% \hyperref[implementation-challenges]{5.9 Implementation Challenges
% \hyperref[implementation-challenges]{35}}

% \hyperref[deployed-streamlit-dashboard-implementation]{5.10 Deployed
% Streamlit Dashboard Implementation
% \hyperref[deployed-streamlit-dashboard-implementation]{36}}

% \hyperref[summary-2]{5.11 Summary \hyperref[summary-2]{36}}

% \hyperref[chapter-6-testing-and-evaluation]{Chapter 6 Testing and
% Evaluation \hyperref[chapter-6-testing-and-evaluation]{37}}

% \hyperref[baseline-results]{6.1 Baseline Results
% \hyperref[baseline-results]{37}}

% \hyperref[proposed-hybrid-model-results]{6.2 Proposed Hybrid Model
% Results \hyperref[proposed-hybrid-model-results]{37}}

% \hyperref[overall-model-comparison]{6.3 Overall Model Comparison
% \hyperref[overall-model-comparison]{38}}

% \hyperref[per-class-performance-of-the-best-model]{6.4 Per-Class
% Performance of the Best Model
% \hyperref[per-class-performance-of-the-best-model]{40}}

% \hyperref[diagnostic-error-analysis]{6.5 Diagnostic Error Analysis
% \hyperref[diagnostic-error-analysis]{42}}

% \hyperref[summary-3]{6.7 Summary \hyperref[summary-3]{42}}

% \hyperref[chapter-7-conclusion-and-future-work]{Chapter 7 Conclusion and
% Future Work \hyperref[chapter-7-conclusion-and-future-work]{43}}

% \hyperref[_Toc239502989]{\textbf{7.1 Conclusion}
% \hyperref[_Toc239502989]{43}}

% \hyperref[_Toc239502990]{\textbf{7.2 Discussion}
% \hyperref[_Toc239502990]{43}}

% \hyperref[_Toc239502991]{\textbf{7.3 Limitations}
% \hyperref[_Toc239502991]{44}}

% \hyperref[_Toc239502992]{\textbf{7.4 Practical CTI Triage Deployment
% Pattern} \hyperref[_Toc239502992]{45}}

% \hyperref[_Toc239502993]{\textbf{7.5 Reflection on Project Development}
% \hyperref[_Toc239502993]{45}}

% \hyperref[_Toc239502994]{\textbf{7.6 Research Contributions}
% \hyperref[_Toc239502994]{45}}

% \hyperref[_Toc239502995]{\textbf{7.7 Project Struggles and Solutions}
% \hyperref[_Toc239502995]{46}}

% \hyperref[_Toc239502996]{\textbf{7.8 Future Work}
% \hyperref[_Toc239502996]{46}}

% \hyperref[_Toc239502997]{\textbf{7.9 Detailed Model Comparison and
% Design Reasoning} \hyperref[_Toc239502997]{46}}

% \hyperref[_Toc239502998]{\textbf{7.10 Feature Importance Discussion}
% \hyperref[_Toc239502998]{46}}

% \hyperref[_Toc239502999]{\textbf{7.11 Additional Analytical Discussion}
% \hyperref[_Toc239502999]{46}}

% \hyperref[_Toc239503000]{\textbf{7.12 Final Validation Against Research
% Questions} \hyperref[_Toc239503000]{47}}

% \hyperref[_Toc239503001]{\textbf{7.13 Final Technical Positioning}
% \hyperref[_Toc239503001]{47}}

% \hyperref[_Toc239503002]{\textbf{7.14 Summary}
% \hyperref[_Toc239503002]{47}}

% \hyperref[_Toc239503003]{\textbf{7.15 Final Word on Technical Depth}
% \hyperref[_Toc239503003]{47}}

% \hyperref[references]{References \hyperref[references]{48}}

% \hyperref[appendix-a]{Appendix A \hyperref[appendix-a]{51}}

% \hyperref[a.1-streamlit-dashboard-screenshots]{A.1 Streamlit Dashboard
% Screenshots \hyperref[a.1-streamlit-dashboard-screenshots]{51}}

% \section{List of Figures}\label{list-of-figures}

% \hyperref[_Toc239502814]{1 The methodology workflow for the CTI
% classification system \hyperref[_Toc239502814]{16}}

% \hyperref[_Toc239502815]{2 Class distribution
% \hyperref[_Toc239502815]{17}}

% \hyperref[_Toc239502816]{3 category distribution donut
% \hyperref[_Toc239502816]{18}}

% \hyperref[_Toc239502817]{4 Word Cloud \hyperref[_Toc239502817]{20}}

% \hyperref[_Toc239502818]{5 text\_length\_outliers
% \hyperref[_Toc239502818]{21}}

% \hyperref[_Toc239502819]{6 entity\_type\_frequency
% \hyperref[_Toc239502819]{21}}

% \hyperref[_Toc239502820]{7 eda\_features\_by\_category
% \hyperref[_Toc239502820]{22}}

% \hyperref[_Toc239502821]{8 feature\_correlation
% \hyperref[_Toc239502821]{22}}

% \hyperref[_Toc239502822]{9 density\_vs\_context\_scatter
% \hyperref[_Toc239502822]{23}}

% \hyperref[_Toc239502823]{10 System architecture of the cybersecurity
% threat classification pipeline \hyperref[_Toc239502823]{28}}

% \hyperref[_Toc239502824]{11 data flow diagram of the cybersecurity
% threat classification pipeline \hyperref[_Toc239502824]{32}}

% \hyperref[_Toc239502825]{12 Models comparison
% \hyperref[_Toc239502825]{39}}

% \hyperref[_Toc239502826]{13 all\_model\_metrics\_grouped\_bar
% \hyperref[_Toc239502826]{39}}

% \hyperref[_Toc239502827]{14 per class metrics visualized
% \hyperref[_Toc239502827]{41}}

% \hyperref[_Toc239502828]{15 User Login and Authentication Panel CYBER
% threat classifier Streamlit App \hyperref[_Toc239502828]{51}}

% \hyperref[_Toc239502829]{16 Enter input CTI text in streamlit dashboard
% \hyperref[_Toc239502829]{52}}

% \hyperref[_Toc239502830]{17 Prediction result of classified threat
% \hyperref[_Toc239502830]{53}}

% \section{List of Tables}\label{list-of-tables}

% \hyperref[_Toc239502831]{1 Comparative Analysis of Recent Literature and
% Proposed System \hyperref[_Toc239502831]{14}}

% \hyperref[_Toc239502832]{2 Baseline Model Results (5-fold
% cross-validation) \hyperref[_Toc239502832]{37}}

% \hyperref[_Toc239502833]{3 Full Model Comparison, Sorted by Macro-F1
% \hyperref[_Toc239502833]{38}}

% \hyperref[_Toc239502834]{4 per class performance of proposed hybrid
% model \hyperref[_Toc239502834]{40}}
\cleardoublepage
\pagenumbering{arabic}
\setcounter{page}{1}
\chapter*{Abstract}
\addcontentsline{toc}{chapter}{Abstract}


Cyber Threat Intelligence (CTI) reporting has grown from simple
indicator lists into narrative-driven text describing adversary tactics,
techniques, and procedures, threat actors, and targeted infrastructure.
Automatically classifying this text into actionable threat categories
has become an important applied challenge, yet the datasets available
for this task rarely include a ready-made category label, and published
hybrid Transformer-recurrent architectures for this task consistently
rely on general-purpose encoders despite domain-adapted cybersecurity
encoders being shown, separately, to represent security vocabulary more
effectively. This dissertation addresses both gaps directly. A genuine
classification target is constructed from raw CTI text using
keyword-rule matching supported by curated malware, ransomware, and
threat-actor lexicons, deliberately avoiding any dependency on
annotator-authored explanatory fields so the label remains learnable
from exactly the information available to a deployed system. A hybrid
SecureBERT-BiLSTM-Attention architecture is developed and benchmarked
against classical machine learning baselines (TF-IDF with logistic
regression and a support vector machine), a bidirectional LSTM trained
from randomly initialised embeddings, and a published literature
benchmark for an equivalent general-BERT hybrid architecture. The
proposed model achieves a macro-averaged F1 score of 0.9370, exceeding
the strongest classical baseline by 2.73 points and the published
general-BERT hybrid benchmark by 7.7 points, directly confirming that
domain adaptation compounds with the architectural benefit of a hybrid
design. A diagnostic evaluation further links model performance to
entity-aware structural measures computed from the
dataset\textquotesingle s existing entity annotations, identifying
low-signal, structurally ambiguous records as the consistent source of
remaining classification difficulty across every model evaluated.
\chapter{Introduction}
% \section{Chapter 1 Introduction}\label{chapter-1-introduction}

\protect\phantomsection\label{_Toc239502922}{}\section{Overview}

Cyber Threat Intelligence (CTI) has become a central input to modern
security operations, providing the situational detail analysts need to
identify, prioritise, and respond to emerging threats. As CTI reporting
has matured beyond simple indicators of compromise into narrative
descriptions of adversary behaviour, automatically classifying this text
into meaningful threat categories has emerged as a practical, high-value
application of natural language processing within cybersecurity. This
dissertation investigates whether a hybrid deep learning architecture,
combining a cybersecurity domain-adapted transformer encoder with a
recurrent sequence layer, can classify CTI text more effectively than
existing approaches, and whether the entity annotations already present
in many CTI datasets can be used analytically to explain where such a
system succeeds and where it does not. The remainder of this chapter
sets out the background, rationale, problem, aims, and structure of the
work that follows.

\protect\phantomsection\label{_Toc239502923}{}\section{Background}

Natural language processing has been applied to cybersecurity text for
over a decade, beginning with rule-based and information-extraction
approaches to pulling structured indicators out of unstructured reports,
and progressing, alongside the wider field, toward transformer-based
language models. Two developments are particularly relevant to this
dissertation. First, domain-adapted transformer encoders, pretrained
specifically on cybersecurity-domain text rather than general internet
text, have been shown to represent security-specific vocabulary, malware
family names, CVE identifiers, and threat-actor terminology, more
effectively than general-purpose encoders such as standard
BERT{[}\href{file:///D:/Downloads/final-report-formatted-v4(1).docx\#ref_3}{\ul{3}},\href{file:///D:/Downloads/final-report-formatted-v4(1).docx\#ref_4}{\ul{4}}{]}.Second,
a growing line of applied work has shown that hybrid architectures,
combining a transformer encoder with a recurrent layer such as a
bidirectional LSTM, outperform either component used alone on
threat-text classification
tasks{[}\href{file:///D:/Downloads/final-report-formatted-v4(1).docx\#ref_10}{\ul{10}}{]}.These
two developments have, until now, been investigated separately: hybrid
architectures in the literature use general-purpose encoders, and
domain-adapted encoders have not been tested inside a hybrid
architecture. CTI datasets have also grown richer in structure,
frequently including annotated entities (malware, threat-actor,
vulnerability, and related types) and relations between them, alongside
the underlying text, following conventions similar to the Structured
Threat Information Expression (STIX) object model widely used in the
threat-intelligence-sharing community.

\protect\phantomsection\label{_Toc239502924}{}\section{Rationale}

Three observations motivate the specific approach taken in this
dissertation. First, since hybrid transformer-recurrent architectures
and domain-adapted encoders have each independently been shown to
improve on a general-purpose, non-hybrid baseline, combining them is a
natural next step that, at the time of writing, had not been tested and
reported in the published literature for CTI classification
specifically. Second, entity and relation annotations, though present in
many CTI datasets at no additional collection cost, are typically used
only to train named-entity recognition models and are rarely turned back
on the dataset itself as a descriptive or diagnostic tool, an
opportunity this dissertation treats as worth pursuing rather than as an
afterthought. Third, a recent systematic review of over forty papers on
threat-technique extraction from CTI text found results largely
incomparable across studies, due to inconsistent ontologies, non-public
datasets, and differing evaluation
protocols{[}\href{file:///D:/Downloads/final-report-formatted-v4(1).docx\#ref_2}{\ul{2}}{]},
motivating the transparent, directly reproducible benchmarking protocol
adopted throughout this work.

\protect\phantomsection\label{_Toc239502925}{}\section{Problem
Statement}

The dataset used in this dissertation, in common with many CTI
resources, does not contain a usable document-level threat-category
label: the closest candidate field is a free-text diagnostic explanation
rather than a fixed category, and the annotated entity fields describe
specific mentions within a record rather than that
record\textquotesingle s overall topic. A genuine classification target
therefore has to be constructed before any classification model can be
trained or evaluated, and constructed in a way that does not depend on
information a deployed system, receiving only raw incoming text, would
not have available. Separately, the dataset is imbalanced and unevenly
distributed across threat categories, properties shown during this
project to measurably disadvantage models, such as a bidirectional LSTM,
that must learn their representations entirely from the available
training data rather than from prior pretraining. Finally, as set out in
Section 1.2, no published hybrid Transformer-BiLSTM architecture for CTI
or related threat-text classification has been evaluated with a
domain-adapted encoder in place of a general-purpose one, leaving open
whether the two independently-demonstrated sources of improvement,
domain adaptation and hybrid architecture, compound, are redundant, or
interact in some other way when combined.

\protect\phantomsection\label{_Toc239502926}{}\section{Research Aim}

The aim of this dissertation is to design, implement, and rigorously
evaluate a hybrid Securebert-BiLSTM-Attention architecture that
substitutes a domain-adapted cybersecurity encoder for a general-purpose
one, benchmarked transparently against classical and recurrent baselines
and against the published literature.

\protect\phantomsection\label{_Toc239502927}{}\section{Research
Objectives}

To achieve the central aim of this study, the following technical
objectives are pursued:\\
1. To construct and validate a genuine, text-derivable threat-category
label for a CTI dataset that does not natively contain one.

2. To implement and benchmark classical machine learning and recurrent
baseline classifiers under a consistent, stratified cross-validation
protocol.

3. To perform a structured, entity-aware exploratory analysis of the
dataset, computing descriptive measures (Entity Density, Entity
Richness, Threat Context Score, Rare Entity Ratio) directly from
existing entity annotations.

4. To design, implement, and evaluate a hybrid
SecureBERT-BiLSTM-Attention architecture, and to benchmark it against
the baselines above and against a comparable result reported in the
published
literature{[}\href{file:///D:/Downloads/final-report-formatted-v4(1).docx\#ref_10}{\ul{10}}{]}.

\protect\phantomsection\label{_Toc239502928}{}\section{Research
Questions}

Primary research question:

Does substituting a domain-adapted cybersecurity encoder for a
general-purpose transformer encoder, within an otherwise identical
hybrid Securebert-BiLSTM-Attention architecture, improve CTI threat
classification performance beyond classical machine learning baselines,
a from-scratch recurrent baseline, and the published literature
benchmark for the equivalent general-purpose-encoder
configuration{[}\href{file:///D:/Downloads/final-report-formatted-v4(1).docx\#ref_10}{\ul{10}}{]}?Supporting
questions:

• Can a genuine, text-derivable threat-category label be constructed for
a CTI dataset that contains no native category field, and validated to
an acceptable standard of accuracy?

• Does resolving preprocessing challenges and class imbalance through
standard pipeline designs measurably affect baseline classification
performance?

\protect\phantomsection\label{_Toc239502929}{}\section{Detailed
Technical Justification}

The choice of SecureBERT as the domain-adapted encoder is justified by
prior published evidence that it outperforms general-purpose RoBERTa and
BERT checkpoints specifically on cybersecurity-domain language modelling
and downstream
tasks{[}\href{file:///D:/Downloads/final-report-formatted-v4(1).docx\#ref_3}{\ul{3}}{]},
making it a reasonable representative of the domain-adaptation approach
rather than an arbitrary selection.The bidirectional LSTM layer
positioned after the encoder is justified on the basis that it re-models
sequential dependency over the encoder\textquotesingle s
already-contextualised token embeddings, a combination with precedent in
adjacent domains such as biomedical and legal named-entity
recognition{[}\href{file:///D:/Downloads/final-report-formatted-v4(1).docx\#ref_9}{\ul{9}}{]},
rather than being a novel or untested architectural pattern.The additive
attention layer is justified both by its function, producing a learned,
weighted summary of the sequence rather than relying on a fixed pooling
operation, and by the interpretability it affords as a secondary
benefit, since attention weights indicate which tokens contributed most
to a given classification decision. Freezing the lower encoder layers
during fine-tuning is justified directly by an empirical finding from
this project\textquotesingle s own baseline results, namely that a
bidirectional LSTM trained entirely from randomly initialised weights
underperforms substantially on a dataset of this scale, indicating that
components required to learn extensively from scratch are disadvantaged
at this sample size; freezing reduces the number of parameters the
hybrid model must similarly learn from scratch. Stratified five-fold
cross-validation, applied uniformly across every model in the
comparison, is justified by the combination of a moderate dataset size
and class imbalance, under which a single train-test split would be
highly sensitive to which records happened to fall into each partition.

\protect\phantomsection\label{_Toc239502930}{}\section{Extended
Technical Discussion}

Beyond the core architectural choices, several further methodological
decisions merit discussion. The threat-category label is constructed
using text alone, not the dataset\textquotesingle s diagnostic
explanation field, a decision made after directly measuring that a
substantial proportion of labels would otherwise depend on information
absent from the raw text a deployed system would receive, a form of
validity gap rather than a purely technical concern. The evaluation
protocol benchmarks the proposed model against a figure reported in the
published literature for a comparable general-purpose-encoder
hybrid{[}\href{file:///D:/Downloads/final-report-formatted-v4(1).docx\#ref_10}{\ul{10}}{]},
rather than against a second, independently retrained general-encoder
hybrid on this project\textquotesingle s own dataset; this is a
deliberate scope decision, made to keep the computational cost of the
evaluation proportionate to the project\textquotesingle s timeframe, and
is treated as a limitation to be stated explicitly rather than
concealed.

\protect\phantomsection\label{_Toc239502931}{}\section{Research
Significance}

This dissertation contributes to the literature by directly testing a
combination, domain-adapted encoder within a hybrid
Transformer-recurrent architecture, that prior published work had
investigated only separately, and by reporting the result transparently
against both an internal baseline suite and an external literature
benchmark{[}\href{file:///D:/Downloads/final-report-formatted-v4(1).docx\#ref_10}{\ul{10}}{]}.It
contributes methodologically by demonstrating a validated approach to
constructing a classification label where a dataset provides none
natively, including the discovery and correction of a label-construction
dependency that would otherwise have compromised the resulting
model\textquotesingle s real-world validity. It contributes an
evaluation approach that goes beyond a single aggregate performance
figure, linking entity-aware structural measures to per-model error
patterns in a way intended to be reusable on other CTI datasets that
share the same entity-annotation structure. Practically, a classifier of
this kind supports the triage stage of security operations, where the
volume of incoming threat reporting routinely exceeds what analysts can
review manually, and a reliable, well-characterised automated first
pass, including an explicit account of which kinds of report it is least
reliable on, has direct operational value.

\protect\phantomsection\label{_Toc239502932}{}\section{Research
Structure}

The remainder of this dissertation is organised as follows. Chapter 2
reviews the literature on natural language processing for cybersecurity,
domain-adapted language models, hybrid transformer-recurrent
architectures, and CTI extraction and classification more broadly,
identifying the specific gap this dissertation addresses. Chapter 3 sets
out the methodology, covering dataset preparation, label construction
and validation, feature engineering, the experimental design and
evaluation protocol, and the baseline and proposed model architectures.
Chapter 4 describes the system design and architecture of the
implemented pipeline in detail. Chapter 5 documents the implementation
itself, including the technology stack, key implementation decisions,
and challenges encountered and resolved during development. Chapter 6
presents the results, benchmarking the proposed hybrid model against the
baselines and the literature, together with per-class and diagnostic
analyses. Chapter 7 discusses these results in relation to the research
questions set out in Section 1.7, alongside the limitations of the work,
and outlines future work and conclusions.

\chapter{Literature Review}


\section{Introduction}\label{introduction}

Cyber Threat Intelligence (CTI) has developed from a collection of
simple Indicators of Compromise (IOCs) to a narrative-driven reporting
of adversary Tactics, Techniques and Procedures (TTPs), threat actors,
targeted infrastructure and recommended defensive
actions\cite{husari2017ttpdrill},\hyperref[ref_2]{\ul{2}}{]}. The
volume and complexity of today\textquotesingle s CTI text has
outstripped the ability of manual analyst review, and Security
Operations Centres (SOCs) and Computer Security Incident Response Teams
(CSIRTs) rely on timely triage of this reporting to prioritise defensive
action. The classification of CTI narratives into actionable threat
categories has thus emerged as an urgent applied challenge in both
academic and operational security research.Modern CTI datasets often
include more than just a narrative text: each record is enriched with
structured entities (e.g., attack-pattern, malware, threat-actor,
campaign, identity, infrastructure, vulnerability, tools, URL, and
location) and explicit relation tuples between these entities. This
annotation style is similar to the Structured Threat Information
Expression (STIX) object model that is widely adopted by the
threat-intelligence-sharing community, and to recent academic resources
like AnnoCTR, which connects entities, tactics, and techniques in real
CTI reports. These annotations are a valuable and underutilized source
of signal. They are typically used for named entity recognition alone,
but are not often used to describe the composition of a dataset or to
explain the success or failure of a classifier on a particular record.
This dissertation does not view that gap as an afterthought, but rather
an opportunity, and entity and relation structure as an analytical lens,
not as raw input to a classifier. Meanwhile, there are tensions on the
modelling side of CTI automation as well. Security text is
under-performed by general-purpose transformer encoders due to the
presence of specialised vocabulary (CVE identifiers, malware family
names, and TTP-specific phrasing), which motivated a family of
domain-adapted encoders like SecureBERT\cite{securebert2022}and
CySecBERT\cite{cysecbert2022}.As the adoption of Large
Language Models (LLMs) for CTI tasks has accelerated, there has been
increasing evidence of unreliable attribution and hallucination when
using LLM models without grounding\cite{ctibench2024}. In
between these two extremes lies a design space of hybrid, well-balanced
architectures that leverage the representation power of transformer
models with sequence modelling and attention, but are still easily
explainable and auditable. This dissertation fills that gap.

\section{Foundational NLP
Architectures}\label{foundational-nlp-architectures}

Modern text classification in cybersecurity, as elsewhere in NLP, rests
on the transformer architecture introduced by Vaswani et
al.\cite{vaswani2017attention}. Unlike recurrent models, the
transformer processes an entire sequence in parallel and models
dependencies between tokens through self-attention, in which each token
representation is updated as a weighted combination of all other tokens
in the sequence.

where Q, K, and V are learned linear projections (query, key, value) of
the input embeddings and d\_k is the key dimension, with the scaling
factor √d\_k preventing the dot products from growing too large and
saturating the softmax. Multiple attention "heads" are computed in
parallel and concatenated, allowing the model to attend to different
types of relationship (e.g., syntactic proximity versus long-range
semantic association) simultaneously. Because self-attention has no
inherent notion of token order, positional encodings are added to the
input embeddings to inject sequence position information.
BERT\cite{bert2018}built on this architecture by
pretraining a bidirectional transformer encoder using masked language
modelling (predicting randomly masked tokens from both left and right
context) together with a next-sentence-prediction objective, producing
contextual embeddings that transfer well to downstream classification
tasks via fine-tuning. BERT and its derivatives now underlie the
majority of modern text-classification systems, including those applied
to cybersecurity text, largely because bidirectional context is well
suited to disambiguating domain terms that depend on surrounding text
(for example, distinguishing a CVE identifier used descriptively from
one used as an indicator of active exploitation).Recurrent architectures
remain relevant despite the dominance of transformers, particularly in
hybrid pipelines. The Long Short-Term Memory network
(LSTM;\cite{lstm1997}) addresses the vanishing-gradient
problem of simple recurrent networks through a gated cell state, with
input, forget, and output gates controlling what information is written
to, retained in, and read from the cell at each
timestep.{[}\hyperref[ref_13]{13}{]} Supports the use of LSTM with
historical load and weather/time-related variables.A Bidirectional LSTM
(BiLSTM) runs two such LSTMs over a sequence (one left-to-right and one
right-to-left) and concatenates their hidden states, giving each
position access to both preceding and following context in a manner that
complements, rather than duplicates, what a transformer encoder already
\cite{cybert2021}
BiLSTM layer exists in adjacent domains: Jiang et
al.\cite{jiang2019bertbilstmcrf}propose a BERT-BiLSTM-CRF architecture
for Chinese electronic medical record named-entity recognition, in which
BERT contextual embeddings are passed through a BiLSTM to model
sequential label dependencies before a Conditional Random Field layer
performs structured prediction. While the present
dissertation\textquotesingle s task is sequence-level classification
rather than sequence labelling (and therefore does not require a CRF
layer), the underlying rationale (that a BiLSTM layer can usefully
remodel sequential dependency over already-contextualised transformer
embeddings) directly motivates the hybrid architecture proposed in
Chapter 4.It is worth noting a practical limitation of the
self-attention mechanism that is directly relevant to CTI text:
computing pairwise attention scores across all token pairs incurs O(n²)
time and memory complexity in sequence length n, which constrains the
maximum practical input length for transformer encoders applied to
long-form CTI narratives without truncation or chunking. CTI reports
vary considerably in length (from short indicator summaries to
multi-paragraph incident writeups) so this quadratic cost is a genuine,
if manageable, engineering constraint on the encoder design discussed in
Chapter 4, rather than a purely theoretical concern. Earlier,
non-transformer embedding approaches such as static word vectors and the
contextual, character-aware embeddings introduced by ELMo predate BERT
and do not suffer this quadratic cost, but they lack the deep
bidirectional contextualization that later made transformer encoders the
default choice for classification tasks; they are noted here for
historical completeness rather than as candidates for the present
architecture.

\section{Domain-Adapted Language Models for Cybersecurity}

General-purpose language models underperform on cybersecurity text for a
well-documented reason: their vocabularies and pretraining corpora
contain little of the specialised terminology that dominates security
narratives, including CVE and CWE identifiers, malware family names,
protocol and tool names, and TTP-specific phrasing drawn from frameworks
such as MITRE ATT\&CK. Subword tokenisers trained on general corpora
tend to fragment such terms into uninformative pieces, degrading the
quality of the resulting contextual embeddings. This limitation has
motivated a line of domain-adapted encoders trained or further
pretrained specifically on cybersecurity corpora. SecureBERT
\cite{securebert2022}
adapts the RoBERTa architecture using a custom tokeniser and continued
pretraining on a large corpus of cybersecurity text, improving
performance on downstream tasks such as named entity recognition and
text classification relative to general BERT/RoBERTa checkpoints.
CySecBERT
\cite{cysecbert2022}
pursues a similar domain-adaptation strategy independently, further
pretraining BERT on a large, curated cybersecurity corpus and reporting
consistent gains over general-domain BERT across several cybersecurity
NLP benchmarks. Across all efforts, the consistent finding is that
domain-adapted encoders outperform general-purpose transformers on
downstream CTI tasks, which directly motivates this
dissertation\textquotesingle s choice to use a domain-adapted
transformer (rather than a general-domain one) as the encoder backbone
for both a standalone baseline and the proposed hybrid model. A common
thread across these domain-adaptation efforts is the treatment of
tokenisation itself as a design decision rather than a fixed
preprocessing step. General-purpose subword tokenisers, trained on web
and encyclopaedic text, tend to split security-specific tokens (a CVE
identifier, a malware family name, a living-off-the-land binary name)
into fragments that carry little standalone meaning, forcing the encoder
to reconstruct domain semantics from noisy subword pieces.
Domain-adapted tokenisers, by contrast, are trained or extended on
cybersecurity corpora so that such tokens are more often preserved as
coherent units, which in turn improves the quality of the contextual
embeddings the encoder produces. An early pioneer in this space was
CyBERT
{[}\href{file:///D:/Downloads/literature-review-and-references-updated.docx\#ref_16}{\ul{16}}{]},
which demonstrated the viability of fine-tuning base BERT encoders on
open-source unstructured and semi-unstructured threat reports to extract
domain-specific entities in Security Operations Center (SOC) workflows.
More recently, this domain-adaptation lineage has culminated in
SecureBERT 2.0
\cite{securebert2025},
which transitions from the standard RoBERTa/BERT base to the ModernBERT
architecture. SecureBERT 2.0 introduces advanced long-context modeling
and hierarchical encoding capabilities to process highly complex and
heterogeneous documents (such as source code and raw threat reports)
while being pretrained on a cybersecurity corpus exceeding 13 billion
text tokens and 53 million code tokens, establishing a new
state-of-the-art for semantic search, entity extraction, and
vulnerability detection.
\section{Hybrid Transformer--Recurrent Architectures for Threat Text Classification}

A more recent and directly relevant line of work combines transformer
encoders with recurrent or graph-based layers specifically for
threat-text classification and extraction, rather than using either
family of model in isolation. Hidayatulloh et al.
\cite{hidayatulloh2026hybrid}
report that a hybrid BERT-BiLSTM architecture, in which BERT-derived
token embeddings are passed through a BiLSTM layer before
classification, improves macro F1 from 0.82 for standalone fine-tuned
BERT to 0.86 on a CTI text classification task, outperforming logistic
regression (0.69), SVM (0.72), and standalone LSTM (0.75) baselines
under the same evaluation protocol. Xiang, Shi, and Zhang
\cite{xiang2023apt}
combine BERT with a BiGRU-CRF layer for Advanced Persistent Threat (APT)
event extraction from web text, reporting the highest F1 score among all
compared architectures, including standalone BERT and ERNIE. Liu and Dai
\cite{liu2024sql}
apply a hybrid BERT-LSTM network to SQL injection attack detection,
again outperforming standalone baselines by dynamically extracting
contextual features that neither model captures alone. Ferrag et al.
\cite{ferrag2024llm}
show that a lightweight, purpose-built BERT variant (SecurityBERT)
outperforms CNN- and RNN-based methods, including prior hybrid
GAN-Transformer and CNN-LSTM solutions, on IoT/IIoT threat detection,
achieving 98.2 percent accuracy across fourteen attack types with
sub-150-millisecond inference time. Across all five of these studies,
the encoder component is either a general-purpose BERT checkpoint
fine-tuned directly on the task, or a purpose-built model trained from
scratch; none evaluate an existing, separately validated domain-adapted
cybersecurity encoder, such as SecureBERT
\cite{securebert2022},
as a drop-in replacement within the hybrid architecture. This is a
notable gap given that domain-adapted encoders have been shown,
independently, to represent cybersecurity vocabulary more effectively
than general-purpose transformers. Whether this domain-adaptation
advantage compounds with the architectural advantage of a hybrid
transformer-recurrent design, rather than being redundant with it, has
not previously been tested. This dissertation addresses that question
directly, by replicating the general-BERT hybrid configuration reported
by Hidayatulloh et al.
\cite{hidayatulloh2026hybrid}
as a literature baseline, then substituting SecureBERT for BERT within
an otherwise identical architecture. This design choice is also
supported by recent empirical evaluations. For example, Kaur et al.
\cite{kaur2024comparison}
conducted a comprehensive comparative analysis of standalone LSTM and
Transformer architectures for threat classification, demonstrating that
while LSTM sequence modeling captures short-term local semantics,
transformer attention is vastly superior in modeling long-range
contextual dependencies within multi-stage threat narratives.
Fine-tuning a hybrid transformer-recurrent design---wherein the
sequential recurrent layer refines the deep, multi-head attention
contextual embeddings of a domain-specific encoder---thus combines the
unique strengths of both paradigms while mitigating their individual
limitations.

\section{TTP and Entity/Relation Extraction from CTI Text}

A substantial and longrunning body of work addresses the extraction of
structured information (entities, relations, and TTPs) from unstructured
CTI reports. TTPDrill
\cite{husari2017ttpdrill}
was an early and influential system in this space, combining information
extraction techniques with an ontology mapping threat actions expressed
in free text to a structured threat action representation, enabling
automatic, ontologygrounded extraction of adversarial behaviour from
prose reports. CASIE
\cite{casie2020}
extended this direction to cybersecurity event extraction, identifying
event triggers and arguments (e.g., attack, vulnerability-related, or
ransom events) within CTI text using a neural sequence labelling
approach. The AnnoCTR dataset provides a large, richly annotated
resource linking entities, tactics, and techniques within genuine CTI
reports, offering a valuable benchmark for entity and TTP linking tasks
and a structural precedent for the dataset schema used in this
dissertation. Despite this sustained research activity, a recent
systematic review by Büchel et al.
\cite{buechel2025sok}
(a "Systematisation of Knowledge" (SoK) covering over forty papers on
automated TTP extraction) concludes that reported results are largely
incomparable across studies. The review attributes this to inconsistent
ontologies (different papers ground extraction in different,
incompatible taxonomies), nonpublic or inconsistently partitioned
datasets, and differing evaluation protocols, and it explicitly
questions whether the field can currently demonstrate genuine cumulative
progress. This is a methodological gap that the present dissertation
addresses directly, by adopting transparent, reproducible baselines
evaluated under a single consistent protocol on one dataset, rather than
by contributing yet another incomparable point estimate to the
literature. Comparing these systems along a common set of dimensions is
itself instructive. TTPDrill and CASIE both predate the widespread
availability of large pretrained transformers and rely respectively on
ontology-driven information extraction and neural sequence labelling
over comparatively narrow, taskspecific corpora. AnnoCTR is more recent
and incorporates transformer-based representations, reflecting the same
broader shift toward pretrained encoders documented in Section 2.3, and
is explicitly constructed as a benchmark dataset linking entities,
tactics, and techniques for others to train and evaluate against. None
of the three, however, use the entity or relation annotations they
produce to describe dataset composition beyond basic frequency counts,
and none link structural annotation properties to downstream error
analysis (the specific combination this dissertation pursues). A
further, more specific observation from this strand of literature
motivates the design of the present study: although entity and relation
annotations are common in CTI datasets (precisely because they are
needed to support TTP and event extraction research) they are rarely
used to characterise dataset composition or to explain model behaviour
once extraction or classification has been performed. Existing work
treats entities and relations almost exclusively as extraction targets
or as intermediate representations feeding a downstream extraction
pipeline, not as descriptive or diagnostic signal in their own right.
The difficulty of mapping raw narratives to structured technical
indicators under low-resource and imbalanced conditions remains a major
challenge. To address this, Park and You
\cite{park2024twostage}
proposed a two-stage training methodology to classify CTI sentences into
MITRE ATT\&CK techniques, demonstrating that pretraining on auxiliary
augmented datasets before fine-tuning on primary reports significantly
boosts Macro-F1 scores on low-resource attack patterns. From an
operational gathering perspective, automated threat triage also depends
on efficiently locating relevant CTI documents. In this area,
ThreatCrawl
\cite{threatcrawl2023}
introduced a BERT-based focused crawling framework that dynamically
adapts its crawling path to retrieve highly relevant Open Source
Intelligence (OSINT) and threat articles from the web, illustrating how
document-level BERT classifiers can be integrated into active incident
ingestion pipelines.
\section{Large Language Models in CTI and Reliability
Concerns}\label{large-language-models-in-cti-and-reliability-concerns}

The recent proliferation of general-purpose LLMs has led to their direct
application to CTI tasks, including summarisation, classification, and
TTP identification via prompting. This uptake has, however, exposed a
mismatch between general LLM evaluation and the demands of CTI
reasoning: CTIBench\cite{ctibench2024}was explicitly
constructed because standard LLM benchmarks fail to capture CTI-specific
reasoning tasks such as threat actor attribution, vulnerability
prioritisation, and TTP mapping, and its results highlight that
hallucination and unreliable attribution remain persistent risks when
generative models are used without external grounding (a finding echoed
across the broader LLM-for-security literature).In practical terms, this
means that an LLM asked to classify or summarise a CTI report may
generate a plausible-sounding but factually unsupported threat
attribution, a failure mode that is particularly costly in a security
context where downstream defensive action may depend on the output.
These reliability concerns motivate the present
dissertation\textquotesingle s choice of architecture. Rather than
relying on an ungrounded generative LLM pipeline for classification,
this work adopts a smaller, supervised, well benchmarked hybrid model
whose behaviour can be evaluated exhaustively against held-out labelled
data and whose errors can be diagnostically traced back to measurable
data characteristics (a degree of auditability that prompted, ungrounded
LLM classification does not currently offer).
\section{Identified Research
Gaps}\label{identified-research-gaps}

Synthesising the literature patterns reveals several concrete gaps that
this dissertation is positioned to address:

1. No domain-adapted encoder tested in a hybrid architecture. Published
hybrid transformer-recurrent CTI and threat-text
classifiers{[}\hyperref[ref_10]{\ul{10}},\hyperref[ref_11]{\ul{11}},\hyperref[ref_12]{\ul{12}}{]}all
use general-purpose BERT as the encoder; none test a domain-adapted
cybersecurity encoder such as SecureBERT within the same architecture,
despite domain adaptation being separately shown to outperform
general-purpose transformers on cybersecurity text.

2. Shallow use of entity/relation annotations. CTI datasets that include
entity and relation annotations alongside text rarely use them for
anything beyond named entity recognition itself; they are seldom used to
characterise dataset composition or difficulty.

3. Fragmented, incomparable evaluation. A systematic review of over
forty TTP extraction papers found results largely incomparable across
studies due to custom ontologies and inconsistent
protocols\cite{buechel2025sok}{]}, leaving open whether the field
is genuinely progressing.

4. Unverified LLM reliability. LLM-based CTI systems show persistent
hallucination and attribution errors, motivating benchmarks built
specifically to expose this\cite{ctibench2024}, yet many
applied systems still deploy LLMs without grounding or a simpler,
well-benchmarked alternative.5. Accuracy-only evaluation. Reported
results commonly emphasise overall accuracy while omitting per-class
performance, obscuring weaknesses on less frequent threat categories.
% table idar
\begin{table}[htbp]
\centering
\caption{Comparative Analysis of Recent Literature and Proposed System}
\label{tab:literature-comparison}

\renewcommand{\arraystretch}{1.35}
\small

\begin{tabularx}{\textwidth}{
>{\raggedright\arraybackslash}X
>{\raggedright\arraybackslash}X
>{\raggedright\arraybackslash}X
>{\centering\arraybackslash}m{2.1cm}
>{\centering\arraybackslash}m{2.3cm}
}
\toprule

\textbf{Study/System} &
\textbf{Encoder Used} &
\textbf{Sequence Layer} &
\textbf{Dataset Size} &
\textbf{Entity-Aware Analysis} \\

\midrule

Hidayatulloh et al.~\cite{hidayatulloh2026hybrid} &
General BERT &
BiLSTM &
CTI Narratives &
No \\

Xiang et al.~\cite{xiang2023apt} &
General BERT &
BiGRU-CRF &
Web Threat Data &
No \\

Satyapanich et al.~\cite{casie2020} &
None (Static Vectors) &
BiLSTM-CRF &
CASIE Corpus &
No \\


\textbf{Proposed System} &
\textbf{SecureBERT (Domain-Adapted)} &
\textbf{BiLSTM-Attention} &
\textbf{1,420} &
\textbf{Yes (Full Evaluation)} \\

\bottomrule
\end{tabularx}
\end{table}
The identified gaps map directly to the experimental design developed in
Chapter 3. In particular, the study tests the domain-adaptation gap
through a controlled general-BERT versus SecureBERT hybrid comparison,
addresses shallow use of annotations through entity-aware EDA and
diagnostic analysis, improves comparability through a common stratified
cross-validation protocol, avoids relying on ungrounded generative
classification, and reports macro and per-class results rather than
accuracy alone.
\chapter{Methodology}
\FloatBarrier
\section{Overview}\label{overview}

The study follows a quantitative, experimental and comparative research
design. Each CTI record is assigned one document-level threat category.
The workflow progresses from preprocessing and target construction to
exploratory data analysis (EDA), model training, and evaluation. TF-IDF
with Logistic Regression and SVM provide lexical baselines; BiLSTM
provides a recurrent baseline; and Transformer--BiLSTM--Attention models
provide the hybrid comparison. This progression isolates the
contribution of lexical representation, sequence modelling, contextual
pretraining and cybersecurity-domain adaptation.

\begin{figure}
\centering
\includegraphics[width=4.23864in,height=5.73256in]{images/media/image1.png}
\caption{\protect\phantomsection\label{_Toc239502814}{} The methodology
workflow for the CTI classification system}
\end{figure}
\FloatBarrier
\section{Dataset and Data
Structure}\label{dataset-and-data-structure}

The dataset contains 1,420 records and 17 original columns. Each record
includes an identifier, CTI text, relations, diagnosis, solutions, and
up to three annotated entity slots containing identifiers, labels and
character-level offsets. Available annotations include malware, threat
actors, attack patterns, infrastructure, software, vulnerabilities,
tools, campaigns, identities, URLs, locations, hashes, file paths, time
expressions, IPv4 addresses, registry keys and email addresses. The
dataset therefore combines narrative text with structured cybersecurity
annotations.
\begin{figure}
\centering
\includegraphics[width=6.76744in,height=3.72571in]{images/media/image2.png}
\caption{\protect\phantomsection\label{_Toc239502815}{} Class
distribution}
\end{figure}

The entity and relation annotations are not fused directly into the
principal classifier input. Instead, they are used for exploratory and
diagnostic analysis. This methodological separation is intentional.
Directly feeding entity-derived features into a classifier whose target
is itself constructed from textual threat indicators could make it
difficult to distinguish genuine contextual learning from label-related
shortcuts. The annotations therefore function as an analytical lens for
examining the data and explaining model behaviour.

\section{Data Cleaning and
Preprocessing}\label{data-cleaning-and-preprocessing}

The preprocessing stage checks dimensions, data types, missing values,
duplicate identifiers and structural consistency. Optional entity slots
may be missing; missing identifiers and offsets use sentinel values and
missing labels use NONE. The single missing diagnosis value is handled
safely when text and diagnosis are combined. Text cleaning is
conservative so that CVE identifiers, URLs, hashes and malware names are
retained. Entity labels are stripped and normalised to uppercase, and
the relations field is parsed for structural validation and descriptive
analysis. Entity labels are stripped of surrounding whitespace and
converted to a consistent uppercase representation. This avoids treating
formatting variations such as URL and url as different categories.
Numeric identifier and offset columns are converted to appropriate
numeric types. The relations field is parsed sufficiently to validate
its structure and support relation-level descriptive analysis. These
operations are applied before model construction so that every
experimental stage receives a consistent and reproducible
representation.

\begin{figure}
\centering
\includegraphics[width=5.81395in,height=4.26481in]{images/media/image3.png}
\caption{\protect\phantomsection\label{_Toc239502816}{} category
distribution donut}
\end{figure}
\FloatBarrier
\section{Construction of the Classification
Target}\label{construction-of-the-classification-target}

The dataset does not provide a native document-level threat category, so
a supervised target is constructed from the CTI text and diagnosis
fields. The combined representation is normalised and processed using
priority-based keyword rules. Nine categories are produced: DoS/DDoS
Attack, Ransomware, Phishing / Social Engineering, Vulnerability
Exploitation, Espionage / APT Campaign, Data Breach / Leakage, Malware
Infection, Threat Actor Activity, and Other / Unspecified. Priority
ordering handles overlapping indicators by checking specific categories
before broader ones. A reproducible sample of 50 records is exported for
manual inspection because the target is derived rather than
human-labelled. The rules are not reported as an independent predictive
baseline because evaluating them against labels generated by the same
rules would be circular. Priority ordering is necessary because one
record can contain several overlapping security indicators. Specific
categories are checked before broader categories so that, for example,
ransomware terminology is not automatically absorbed into the general
malware category. Vulnerability and CVE indicators are similarly
considered before broad threat-actor terminology. Records that do not
satisfy any category-specific rule are assigned to Other / Unspecified.
Because the target is derived rather than human-labelled, a reproducible
sample of 50 records is exported for manual inspection. This
quality-control stage is used to assess whether the generated labels
reasonably represent the underlying CTI narrative. The keyword rules are
therefore treated as a target-construction mechanism rather than as an
independent predictive baseline, because evaluating those same rules
against labels generated by them would be circular.
\section{Entity-Aware Exploratory
Features}\label{entity-aware-exploratory-features}

EDA derives record-level measures from the annotations. Entity Count
measures available entities; Entity Density normalises entity count by
text length; and Entity Richness counts unique entity types. Threat
Context Score weights entity concentration, with additional emphasis on
THREAT-ACTOR, CAMPAIGN, VULNERABILITY and ATTACK-PATTERN. Rare Entity
Ratio measures the proportion of entities from relatively infrequent
types. These measures support EDA and error analysis rather than direct
classifier input. A Threat Context Score is calculated as a composite
measure of entity concentration. Each entity contributes a base score,
while high-signal categories such as THREAT-ACTOR, CAMPAIGN,
VULNERABILITY and ATTACK-PATTERN receive additional weighting. The score
is used as a contextual-information proxy rather than being interpreted
as an official security severity rating. Rare Entity Ratio measures the
proportion of entities associated with relatively infrequent entity
types, with the rarity threshold derived from the observed distribution
rather than manually selected. These variables are stored for EDA and
later error analysis rather than used as direct predictors in the main
text-classification models.
\FloatBarrier
\section{Text-Level Statistical
Features}\label{text-level-statistical-features}

Four text statistics are generated: word count, character count, average
word length and lexical diversity. These measures describe record length
and vocabulary variation and support comparison of text complexity
across threat categories and prediction outcomes.

\begin{figure}
\centering
\includegraphics[width=5.97674in,height=3.16258in]{images/media/image4.png}
\caption{\protect\phantomsection\label{_Toc239502817}{} Word Cloud}
\end{figure}


\section{Exploratory Data
Analysis}\label{exploratory-data-analysis}

EDA examines class distribution, text length and entity characteristics
before model comparison. Class plots identify imbalance; histograms and
boxplots show text-length variation; IQR inspection flags unusual
records without automatic deletion; and word-cloud and entity-frequency
plots summarise terminology. Further analysis compares Entity Density,
Threat Context Score and Rare Entity Ratio across categories using
correlation heatmaps and scatter plots. Entity and relation patterns are
inspected descriptively rather than used in a graph neural network.

\begin{figure}
\centering
\includegraphics[width=6.77216in,height=2.38731in]{images/media/image5.png}
\caption{\protect\phantomsection\label{_Toc239502818}{}
text\_length\_outliers}
\end{figure}

\begin{figure}
\centering
\includegraphics[width=6.53488in,height=4.33211in]{images/media/image6.png}
\caption{\protect\phantomsection\label{_Toc239502819}{}
entity\_type\_frequency}
\end{figure}

\begin{figure}
\centering
\includegraphics[width=7.12614in,height=2.34884in]{images/media/image7.png}
\caption{\protect\phantomsection\label{_Toc239502820}{}
eda\_features\_by\_category}
\end{figure}

\begin{figure}
\centering
\includegraphics[width=5.52326in,height=4.98828in]{images/media/image8.png}
\caption{\protect\phantomsection\label{_Toc239502821}{}
feature\_correlation}
\end{figure}

\begin{figure}
\centering
\includegraphics[width=5.46512in,height=3.81287in]{images/media/image9.png}
\caption{\protect\phantomsection\label{_Toc239502822}{}
density\_vs\_context\_scatter}
\end{figure}
\FloatBarrier
Further analysis compares Entity Density, Threat Context Score and Rare
Entity Ratio across categories. A correlation heatmap examines
relationships among the numerical diagnostic variables, while scatter
plots investigate whether entity concentration and contextual complexity
vary by threat class. Entity interaction patterns are also examined
descriptively using the supplied annotations and relations. The purpose
is not to construct a graph neural network, but to determine whether
particular combinations of security entities characterise more complex
records. The EDA outputs subsequently provide the variables used to
stratify diagnostic error analysis.
\section{Model Training and
Validation}\label{model-training-and-validation}

Stratified five-fold cross-validation is used because the constructed
threat categories are imbalanced. Stratification preserves approximately
similar class proportions in each fold and reduces the risk that
minority categories are poorly represented. This is especially important
in cybersecurity classification, where rare threat categories can be
operationally significant even when their sample size is limited. A
fixed seed of 42 supports reproducibility. Out-of-fold predictions are
retained so that each record is predicted by a model that did not train
on that record. This design improves the credibility of evaluation
because it avoids training-test leakage at the record level. The same
fold structure is used across competing models to ensure a fair
comparison. Performance is evaluated with metrics that reflect class
imbalance and multi-class classification behavior. Macro-F1 is treated
as the principal metric because it calculates F1 independently for each
class and then gives equal weight to every class. This prevents majority
categories from dominating the evaluation. Accuracy is less informative
on its own, because a model can appear strong overall while performing
poorly on rare but important threat categories.
\section{Classical Machine Learning
Baselines}\label{classical-machine-learning-baselines}

The classical baseline stage uses TF-IDF unigrams and bigrams with a
maximum vocabulary of 3,000 features. TF-IDF is appropriate because CTI
documents often contain highly informative keywords, entity names, and
attack indicators that can be captured effectively through sparse
lexical representation. Logistic Regression is used as an interpretable
lexical baseline with balanced class weighting and up to 1,000
iterations. This model tests whether linear decision boundaries over
sparse text features are sufficient for the classification task. A
linear Support Vector Machine with balanced class weighting provides a
second benchmark and evaluates whether stronger margin-based separation
improves performance. These models are valuable because they establish a
strong and efficient reference point. If a deep learning model fails to
outperform them, then the additional complexity may not be justified.
Conversely, if contextual models significantly improve macro-F1, the
results would support the argument that CTI classification benefits from
deeper semantic representations.
\section{BiLSTM Deep Learning
Baseline}\label{bilstm-deep-learning-baseline}

A Bidirectional Long Short-Term Memory network is implemented as a deep
learning baseline without pretrained weights. The vocabulary is built
from the dataset, records are mapped to integer sequences, and the input
length is limited and padded to 60 tokens. The network uses a
64-dimensional embedding layer, a bidirectional LSTM with hidden
dimension 64, mean pooling, dropout, and a fully connected
classification layer. Training uses Adam with learning rate 0.001,
weight decay 1e-5, batch size 16, and 15 epochs. This setup tests
whether sequence modeling can capture word-order dependencies and
contextual cues beyond what TF-IDF can represent. BiLSTM is especially
relevant in CTI texts because threat descriptions often depend on local
context, such as whether a term refers to an attacker, a tool, or a
victim. The BiLSTM baseline is also useful as a bridge between classical
machine learning and transformer-based modeling. It evaluates whether a
learnable sequence representation provides a meaningful advantage even
without large-scale external pretraining.
\section{Transformer and Hybrid
Architecture}\label{transformer-and-hybrid-architecture}

The advanced stage is configured around BERT and SecureBERT. Transformer
inputs use a maximum sequence length of 64 with attention masks. The
classifier applies mean pooling over valid token representations
followed by dropout and a linear head. The configured training uses
AdamW, learning rate 2e-5, five epochs, and batch size 8. The first
eight encoder layers are frozen to reduce computational cost and
overfitting risk. The proposed hybrid architecture extends the
transformer encoder with a bidirectional LSTM, attention mechanism, and
classification head. Contextual transformer outputs are processed
sequentially by the BiLSTM, while attention weights emphasize
informative tokens before classification. The hybrid uses an LSTM hidden
dimension of 128. The notebook also provides a general-BERT hybrid for
comparison and identifies the SecureBERT hybrid as the proposed model.

This transformer-based section is important because it introduces
cybersecurity-domain language modelling into the pipeline. SecureBERT is
conceptually attractive for CTI because it is designed to better capture
domain-specific terminology than a general-purpose
encoder\cite{securebert2022}. The hybrid architecture further
strengthens this idea by combining pretrained contextual embeddings with
sequence-level refinement and attention-based focus.

The transformer sections are configured for GPU-enabled Colab and
require Hugging Face dependencies, so the methodology does not assume
numerical results until those cells are executed. This keeps the chapter
method-focused rather than result-dependent.
\section{Evaluation and Error
Analysis}\label{evaluation-and-error-analysis}

Macro-F1 is the principal metric because the task contains multiple
categories and class imbalance. It calculates F1 independently for each
class and then gives equal weight to every class. For each model, mean
macro-F1 and standard deviation across five folds are reported.
Out-of-fold predictions support classification reports and
confusion-matrix analysis for the best-performing model. Diagnostic
error analysis examines prediction correctness across
rare\_entity\_ratio and threat\_context\_score bands. This tests whether
structurally unusual or information-dense CTI records are harder to
classify and links EDA findings to model behavior. Models are ranked by
mean macro-F1, and the highest-performing model is selected as the
experimental candidate for subsequent project stages.
\section{Reproducibility and Experimental
Workflow}\label{reproducibility-and-experimental-workflow}

The workflow is implemented as a sequential Google Colab/Jupyter
Notebook using Pandas, NumPy, Matplotlib, Seaborn, Scikit-learn,
PyTorch, and Hugging Face Transformers for the transformer stage. The
pipeline covers loading, validation, missing-value handling, text
normalisation, entity-label standardisation, relation parsing,
rule-based target construction, feature engineering, EDA, stratified
five-fold validation, baseline training, deep learning, transformer
experimentation, model comparison, and error analysis.
Feature-engineered data and EDA outputs are saved to support
reproducibility. This is important because dissertation work should be
inspectable and repeatable rather than dependent on hidden manual steps.
Reproducibility also supports academic transparency by making the
transformation from raw CTI records to evaluation-ready features
explicit. The methodology clearly separates dataset-provided information
from derived information. Entity annotations are treated as structured
domain knowledge, while the document-level threat category is explicitly
identified as a constructed target. This provides a transparent
experimental basis for comparing lexical, sequential, and
cybersecurity-domain contextual representations while connecting
EDA-driven feature engineering to measurable model behavior.
\section{Methodological Considerations and
Scope}\label{methodological-considerations-and-scope}

Several methodological decisions are intentionally constrained to keep
the experiment reproducible and proportionate to the available dataset.
First, the relations field is parsed only for structural validation and
relation-count inspection; graph neural networks and relation-derived
predictors are outside the present scope. Second, the engineered
threat\_context\_score is treated as a research feature rather than a
validated cybersecurity severity index. Its value is to test whether
domain-informed entity weighting contains predictive information. Third,
the constructed threat category must be interpreted differently from a
manually annotated ground-truth label. The 50-record manual validation
sample therefore acts as an important quality-control step, and the
limitations of rule-derived labels are retained for discussion. Finally,
all comparative models use the same record-level folds so that
performance differences can be attributed more consistently to the
representation and model rather than to different sampling procedures.
These controls strengthen internal experimental validity while keeping
the implementation manageable within the dissertation scope. The
methodology also leaves a clear path for future extensions, including
relation-aware graph features, manually verified labels, domain-specific
modeling, and explainability analysis, without making those components
necessary for the core experimental contribution.
\section{Streamlit Dashboard
Integration}\label{streamlit-dashboard-integration}

To translate the trained predictive models into an operationally useful
tool for security analysts, the final stage of the methodology includes
the integration of an interactive Streamlit web application. This
dashboard abstracts the underlying neural network complexity, providing
a user-friendly interface where a raw security report or CTI text is
given as the input, and the output is a classified threat label along
with real-time confidence scores and defensive recommendations. By
exposing the hybrid SecureBERT-BiLSTM-Attention model through a simple
interactive interface, security operations center (SOC) analysts can
rapidly triage incoming threat intelligence narratives without requiring
deep command-line or machine learning expertise.
\section{Summary}\label{summary}

Overall, the methodology establishes a reproducible pipeline from raw
CTI records to evaluated threat-category predictions. It combines
conservative preprocessing, transparent rule-based label construction,
entity-aware feature engineering, detailed EDA, stratified
cross-validation, classical text classification, BiLSTM sequence
learning, and transformer-based hybrid modelling. The resulting
framework provides both a practical classification system and an
experimental basis for analysing how cybersecurity-specific contextual
information influences model performance.

\chapter{System Analysis and Design}
\section{Overview of the System
Architecture}\label{overview-of-the-system-architecture}

The system is organised as a modular, seven-stage pipeline: data
ingestion, preprocessing and label construction, feature engineering,
exploratory analysis, baseline model training, the proposed hybrid
model, and mitigation recommendation, with evaluation implemented as a
cross-cutting concern rather than a single terminal stage. Each stage
consumes the output of the previous one and can be run, inspected, and
re-validated independently, which was a deliberate design choice given
that several stages, label construction in particular, were revised
during the project once their output was examined directly rather than
assumed correct in advance. The pipeline is implemented as a single,
sequentially executed notebook rather than as separate services, which
is appropriate for a research system evaluated offline rather than one
intended for live deployment, and keeps every intermediate artefact
(cleaned data, constructed labels, engineered features, trained model
predictions) inspectable at the point it is produced.
\begin{figure}
\centering
\includegraphics[width=6.74445in,height=4.48837in]{images/media/image10.png}
\caption{\protect\phantomsection\label{_Toc239502823}{} System architecture of
the cybersecurity threat classification pipeline}
\end{figure}



The dataset comprises 1,420 cyber threat intelligence records, each
containing free text, a diagnostic explanation, recommended solution
text, and up to three annotated entities drawn from an eighteen-type
STIX-style taxonomy, connected via explicit relation tuples. The data
layer exposes a single flat tabular interface to every downstream
module, so that preprocessing, feature engineering, and modelling
components can be developed and tested against a stable schema.
\section{Preprocessing and Label Construction
Module}\label{preprocessing-and-label-construction-module}

This module performs four functions: handling missing values in the
optional second and third entity slots, casting identifier and text
columns to consistent types, normalising text whitespace and
entity-label casing, and constructing the threat\_category target
variable. The label is constructed via priority-ordered keyword-rule
matching over each record\textquotesingle s own text field, combined
with curated lexicons of named malware and ransomware families and named
threat-actor groups, since many records reference a specific family or
group, for example a named ransomware strain, without using a generic
trigger word such as "ransomware" or "campaign". The label is
deliberately derived from text alone, and never from the diagnosis
field, so that it is learnable in principle from exactly the information
the classification models in Sections 5.5 and 5.6 receive as input; this
constraint is enforced architecturally, by simply never passing the
diagnosis field into any component downstream of label construction,
rather than relying on a documentation note alone. A validation
sub-module exports a random sample of records for manual review, treated
as a required check on label quality rather than an optional step.
\section{Feature Engineering
Module}\label{feature-engineering-module}

Two families of descriptive feature are computed directly from existing
columns and persisted back into the dataset: entity-aware measures
(Entity Density, Entity Richness, Threat Context Score, Rare Entity
Ratio) derived from the annotated entity columns, and standard text
statistics (word count, character count, average word length, lexical
diversity) derived from the text column. These are computed once, stored
as columns, and used descriptively in the exploratory analysis module
and diagnostically in the evaluation stage; they are not supplied to any
classification model as an input feature, keeping the classification
architecture itself simple and keeping the analytical and predictive
roles of the entity annotations clearly separated.
\section{Exploratory Analysis
Module}\label{exploratory-analysis-module}

A dedicated visualisation module consumes the outputs of the feature
engineering module to produce class-distribution, text-length-outlier,
word-cloud, entity-frequency, feature-correlation, and per-category
boxplot views, all persisted as image files alongside the
notebook\textquotesingle s inline output. This module has no downstream
dependents in the modelling pipeline; its role is purely to characterise
the dataset and to surface irregularities, such as the file-hash-only
text record identified as a length outlier, before modelling begins.
\section{Baseline Model
Architectures}\label{baseline-model-architectures}

Two non-hybrid baseline families are implemented. The first represents
each record as a TF-IDF vector over unigrams and bigrams, classified
using either logistic regression or a linear support vector machine,
both with class weights set inversely proportional to category
frequency. The second is a bidirectional LSTM operating over trainable,
randomly initialised word embeddings, using mean pooling over the
sequence output before a linear classification head. Both families are
deliberately simple, non-pretrained architectures, providing a lower
bound against which the benefit of pretrained transformer
representations in the proposed architecture can be measured, and both
are wrapped in the same stratified five-fold evaluation harness
described in Section 4.8 so that every model in the comparison is
subject to an identical protocol.
\section{Proposed Hybrid
Architecture}\label{proposed-hybrid-architecture}

The proposed architecture follows a four-layer design:

Layer 1, Encoder: SecureBERT, a transformer encoder pretrained on
cybersecurity-domain text, converts each record\textquotesingle s text
into contextual token embeddings. The lower encoder layers are frozen
during fine-tuning; only the upper layers and the layers described below
are trained, reducing the risk of overfitting a large pretrained model
to a dataset of this size.

Layer 2, Sequential modelling: A bidirectional LSTM re-models sequential
dependency over the encoder\textquotesingle s token embeddings,
producing a sequence of hidden states that combine the
transformer\textquotesingle s contextual representation with an explicit
forward and backward pass over the sequence.\\
Layer 3, Attention: An additive attention mechanism computes a weighted
importance score for each position in the BiLSTM output, producing a
single weighted summary vector per record rather than relying on a fixed
pooling operation such as mean or max pooling, and incidentally
providing a per-token importance score that could support explainability
work beyond the scope of this dissertation.

Layer 4, Classification head: A linear layer maps the attention-weighted
summary vector to the nine threat categories.
\section{Evaluation Harness}\label{evaluation-harness}

A shared training-and-evaluation function underlies every
transformer-based model in the comparison, parameterised by a
model-constructor callback so that a structurally identical training
loop, differing only in which encoder the constructor returns, is used
for both the replicated and proposed hybrid configurations. This
function performs stratified five-fold splitting, per-fold model
instantiation and training, out-of-fold prediction collection, and
per-fold macro-F1 computation, so that the resulting comparison table in
Chapter 6 reflects a single, consistently applied evaluation protocol
across every model rather than several ad hoc ones.
\section{Design Rationale
Summary}\label{design-rationale-summary}

Three architectural decisions recur across this chapter and are worth
stating explicitly. First, every component that produces the
classification label or the classification prediction is kept strictly
separate from components that only describe or diagnose the data, so
that entity-aware and text-statistic features can be developed and
extended freely without risk of silently changing what the models are
trained on. Second, the same evaluation harness is applied uniformly to
every model under comparison, including the two hybrid configurations
that differ only in encoder choice, so that differences in the resulting
comparison table can be attributed to genuine architectural or
representational differences rather than to inconsistencies in how each
model happened to be evaluated. Third, encoder layer freezing and a
bounded fine-tuning budget are treated as architectural choices rather
than incidental hyperparameters, directly motivated by the
BiLSTM-from-scratch result in Chapter 3 showing that this
dataset\textquotesingle s size penalises components trained entirely
without prior knowledge.
\section{Summary}\label{summary-1}

In summary, the system analysis and design defines a rigorous, robust,
and cleanly separated architectural pipeline. This structural discipline
prevents evaluation leakage and ensures the integrity of the core
encoder-comparison experiment.

\chapter{Implementation}
\section{Technology Stack}\label{technology-stack}

The system is implemented in Python, using pandas and numpy for data
handling, scikit-learn for classical baselines and evaluation metrics,
PyTorch and the Hugging Face transformers library for the BiLSTM and
hybrid transformer models, and matplotlib, seaborn, and wordcloud for
visualisation. Transformer fine-tuning is executed on Google Colab using
a GPU runtime, since encoder fine-tuning is impractical on CPU at a
reasonable iteration speed, while data preparation, feature engineering,
classical baselines, and the BiLSTM baseline require no GPU and were
developed and verified independently of the transformer sections. The
full pipeline is implemented as a Jupyter notebook, structured so that
each stage in Chapter 5 corresponds to a distinct, independently
re-runnable section, with intermediate datasets (the cleaned dataset,
the feature-engineered dataset) saved to disk between major stages.

\FloatBarrier
\section{Data Preparation
Implementation}\label{data-preparation-implementation}

\begin{figure}
\centering
\includegraphics[width=6.74445in,height=4.48837in]{images/media/image11.png}
\caption{\protect\phantomsection\label{_Toc239502824}{} data flow
diagram of the cybersecurity threat classification pipeline}
\end{figure}


Missing-value handling was initially implemented for the optional third
entity slot only, matching the pattern observed in the dataset, where
every record had at least two annotated entities. The handling was
extended to cover the second slot identically, a concrete example of an
implementation assumption that held for the original data but did not
generalise once the data changed. A separate, unrelated defect was identified in
the text-cleaning function, where a whitespace-collapsing regular
expression was written with an extra level of escaping, causing it to
silently fail to collapse repeated whitespace without raising an error;
this was corrected once identified by direct inspection of the cleaned
output.
\section{Label Construction
Implementation}\label{label-construction-implementation}

The categorisation function checks eight thematically ordered pattern
groups against each record\textquotesingle s lowercased text, returning
the first matching category, with an explicit "Other/Unspecified"
fallback when no pattern matches. Ordering matters functionally, not
just stylistically: more specific patterns, for example
ransomware-specific phrasing, are checked before more general ones, such
as generic malware phrasing, so that a specific pattern is not shadowed
by a broader one earlier in the sequence. Two curated lexicons, one of
named malware and ransomware families and one of named threat-actor and
APT group names, are checked alongside the regular-expression patterns
within their respective category checks, since manual inspection of the
initial "Other/Unspecified" bucket showed that a substantial share of
unmatched records named a specific family or group, for example a
malware family or an APT group, without using a generic trigger word;
adding these lexicons reduced the Other/Unspecified category from
roughly a fifth of the dataset to under a tenth, without altering the
underlying text-only construction principle.
\section{Feature Engineering
Implementation}\label{feature-engineering-implementation}

Entity-aware features are computed by iterating over each
record\textquotesingle s populated entity-label columns, counting
entities, counting distinct types, and computing a rarity threshold from
the twenty-fifth percentile of global type frequency. The Threat Context
Score is implemented as a weighted sum of entity count and the count of
high-signal entity types (threat-actor, campaign, vulnerability,
attack-pattern), deliberately excluding any relation-derived term,
consistent with the scope decision recorded in Chapter 4. Text
statistics are computed directly from the text column using standard
Python string operations, requiring no additional dependencies, and all
engineered columns are persisted into the dataset once computed, rather
than recomputed inside every downstream cell that uses them.
\section{Baseline Model
Implementation}\label{baseline-model-implementation}

The TF-IDF baselines use scikit-learn\textquotesingle s TfidfVectorizer
with a three-thousand-feature cap and a unigram/bigram range, feeding
logistic regression and a linear-kernel support vector machine, both
evaluated via stratified five-fold cross-validation using
cross\_val\_score for the headline macro-F1 figures and
cross\_val\_predict for out-of-fold predictions, the latter providing
the per-class and diagnostic analyses in Chapter 6 without requiring a
separate held-out split. The BiLSTM is implemented as a small PyTorch
module, an embedding layer, a bidirectional LSTM, mean pooling, and a
linear head, trained from randomly initialised weights independently on
each cross-validation fold for fifteen epochs with the Adam optimiser
and a small weight-decay term, so that no information from a
fold\textquotesingle s held-out data influences that
fold\textquotesingle s training, and so that five independently trained
models, rather than one model evaluated five times, produce the reported
score.
\section{Hybrid Model
Implementation}\label{hybrid-model-implementation}

The hybrid model reuses a shared training-loop function across both the
replicated general-BERT configuration and the proposed SecureBERT
configuration, taking a model-constructor function as an argument so
that a fresh, untrained model is instantiated for every cross-validation
fold rather than a single model being reused and implicitly leaking
information across folds. Within each fold, encoder parameters in the
lower eight transformer layers are frozen, excluded from the optimiser
via a parameter filter, and the remaining parameters are fine-tuned
using AdamW with a small learning rate appropriate for transformer
fine-tuning, for a small number of epochs given the dataset size.
Batches are tokenised with padding and truncation to a fixed maximum
sequence length, and attention masks are used both by the encoder itself
and by the additive attention layer, to ensure padded positions do not
contribute to the attention-weighted summary. Each
fold\textquotesingle s model and optimiser are explicitly deleted and
the GPU cache cleared before the next fold begins, since retaining five
full transformer models simultaneously in GPU memory is unnecessary and
would otherwise risk exhausting available memory on a standard Colab GPU
allocation.
\section{Evaluation and Reporting
Implementation}\label{evaluation-and-reporting-implementation}

A full metrics report is computed for every model by extracting
accuracy, macro-precision, macro-recall, macro-F1, and, where predicted
probabilities are available, macro-averaged one-vs-rest ROC-AUC,
assembled into a single comparison table sorted by macro-F1 to identify
the best-performing model automatically rather than by manual
inspection. For the best model specifically, per-class precision,
recall, and F1 are reported alongside a normalised confusion matrix and,
where probabilities are available, ROC curves per class and a
prediction-confidence distribution split by correctness, so that the
evaluation extends beyond a single aggregate score into a per-class and
per-prediction view of model behaviour.
\section{Reproducibility
Measures}\label{reproducibility-measures}

A fixed random seed is set once at the start of the pipeline and reused
across numpy, PyTorch, and every scikit-learn estimator that accepts a
random\_state argument, and the same stratified five-fold split object
is reused across every model in the comparison, so that all models are
evaluated on identical fold partitions rather than independently
resampled ones. Intermediate artefacts, the cleaned and relabelled
dataset, the feature-engineered dataset, the fifty-record
label-validation sample, and every generated figure, are written to disk
at the point they are produced, so that any stage of the pipeline can be
re-inspected or resumed without re-running the stages before it.
\section{Implementation
Challenges}\label{implementation-challenges}

Three implementation issues were identified and resolved during
development, each retained here as a documented part of the process
rather than silently corrected without record:

1. A text-cleaning regular expression that failed to collapse whitespace
due to over-escaping.

2. A missing-value handling routine that did not generalise properly
when processing optional entity fields with missing labels.

3. A label-construction dependency on the diagnosis field that made
roughly a quarter of constructed labels unreachable from the text a
deployed system would actually receive. This was resolved by
reconstructing the label from text alone, expanding the
keyword-and-lexicon rules to limit the resulting growth of the
Other/Unspecified category, and validating the resulting category
distribution before it was used in any result reported in Chapter 6..
\section{Deployed Streamlit Dashboard
Implementation}\label{deployed-streamlit-dashboard-implementation}

Following the offline training and evaluation of the hybrid
SecureBERT-BiLSTM-Attention model, a web-based interactive triage tool
was implemented using the Streamlit framework. The application loads the
serialized model weights, preprocessor, and tokenizer into memory.
Within the dashboard, security analysts can input unstructured text,
such as a newly received threat report or campaign write-up. The
application pre-processes this text run-time, extracts token
representations via the SecureBERT encoder, passes them through the
BiLSTM and attention layers, and outputs a classified threat label in
real-time. This interactive interface successfully bridges offline
machine learning research with practical, low-latency triage operations,
allowing immediate classification of security narratives.
\section{Summary}\label{summary-2}

In summary, the implementation chapter details the practical translation
of architectural designs into executable, modular PyTorch and
scikit-learn pipelines, ensuring high-fidelity reproducibility of the
entire threat-intelligence classification model suite.
\chapter{Testing and Evaluation}
\section{Baseline Results}\label{baseline-results}

All baseline models were evaluated using stratified five-fold
cross-validation on the full 1,420-record dataset, with macro-averaged
F1 as the primary metric, given the class imbalance across the nine
threat categories described in Chapter 3.

\begin{table}[htbp]
\centering
\caption{Baseline Model Results (5-fold Cross-Validation)}
\label{tab:baseline-results}

\renewcommand{\arraystretch}{1.35}
\small

\begin{tabularx}{\textwidth}{
>{\raggedright\arraybackslash}X
>{\centering\arraybackslash}m{2.1cm}
>{\centering\arraybackslash}m{2.2cm}
>{\centering\arraybackslash}m{2.2cm}
>{\centering\arraybackslash}m{2cm}
}
\toprule

\textbf{Model} &
\textbf{Accuracy} &
\textbf{Precision (Macro)} &
\textbf{Recall (Macro)} &
\textbf{F1 (Macro)} \\

\midrule

TF-IDF + Logistic Regression & 0.8937 & 0.9232 & 0.8933 & 0.9054 \\

TF-IDF + SVM & 0.8993 & 0.9262 & 0.8979 & 0.9097 \\

BiLSTM (from scratch) & 0.8507 & 0.8716 & 0.8568 & 0.8632 \\

\bottomrule
\end{tabularx}
\end{table}
According to the above results \textbf{table} the TF-IDF and support
vector machine baseline achieved the highest macro-F1 among the
non-hybrid models, narrowly ahead of logistic regression, with the
BiLSTM trailing both, consistent with the finding that sequence-level
models trained from scratch without large-scale pretraining are
sensitive to the training set size.
\section{Proposed Hybrid Model
Results}\label{proposed-hybrid-model-results}

The proposed hybrid SecureBERT-BiLSTM-Attention model achieved a
macro-F1 of 0.9372 (standard deviation 0.0132), with an accuracy of
0.9359, a macro-precision of 0.9422, and a macro-recall of 0.9336,
exceeding every baseline in Section 6.1, including the strongest
classical baseline (TF-IDF and SVM, macro-F1 0.9097), by 2.73 macro-F1
points. For reference, Hidayatulloh et
al.{[}\href{file:///D:/Downloads/final-report-formatted-v4(1).docx\#ref_10}{\ul{10}}{]}
report a macro-F1 of 0.86 for an equivalent hybrid architecture using a
general-purpose BERT encoder on a comparable CTI classification task;
the proposed SecureBERT-based configuration exceeds this published
figure by 7.7 macro-F1 points. While the datasets, label taxonomies, and
evaluation protocols are not identical this margin is large enough to be
a meaningful, rather than marginal, result in favour of the
domain-adapted encoder.
\section{Overall Model
Comparison}\label{overall-model-comparison}
\begin{table}[htbp]
\centering
\caption{Full Model Comparison, Sorted by Macro-F1}
\label{tab:full-model-comparison}

\renewcommand{\arraystretch}{1.35}
\small

\begin{tabularx}{\textwidth}{
>{\raggedright\arraybackslash}X
>{\centering\arraybackslash}m{2cm}
>{\centering\arraybackslash}m{2.2cm}
>{\centering\arraybackslash}m{2.2cm}
>{\centering\arraybackslash}m{2cm}
}
\toprule

\textbf{Model} &
\textbf{Accuracy} &
\textbf{Precision (Macro)} &
\textbf{Recall (Macro)} &
\textbf{F1 (Macro)} \\

\midrule

BiLSTM (from scratch) & 0.8507 & 0.8716 & 0.8568 & 0.8632 \\

TF-IDF + Logistic Regression & 0.8937 & 0.9232 & 0.8933 & 0.9054 \\

TF-IDF + SVM & 0.8993 & 0.9262 & 0.8979 & 0.9097 \\

 
\textbf{Hybrid SecureBERT-BiLSTM-Attention (Proposed)} &
\textbf{0.9359} &
\textbf{0.9422} &
\textbf{0.9336} &
\textbf{0.9372} \\

\bottomrule
\end{tabularx}
\end{table}
According to above \textbf{table} the proposed hybrid model is the
best-performing configuration overall.The proposed hybrid
SecureBERT-BiLSTM-Attention model achieved a macro-F1 of 0.9372
(standard deviation 0.0132), with an accuracy of 0.9359, a
macro-precision of 0.9422, and a macro-recall of 0.9336, exceeding every
baseline , confirming the central hypothesis of this dissertation:
substituting a domain-adapted encoder into an otherwise identical hybrid
Transformer-BiLSTM-Attention architecture improves performance beyond
what the same architecture achieves with a general-purpose encoder in
the literature, and beyond every non-hybrid baseline evaluated on this
dataset

\begin{figure}
\centering
\includegraphics[width=6.5in,height=3.57847in]{images/media/image12.png}
\caption{\protect\phantomsection\label{_Toc239502825}{} Models
comparison}
\end{figure}
\FloatBarrier
\begin{figure}
\centering
\includegraphics[width=3.85417in,height=3.53403in]{images/media/image13.png}
\caption{\protect\phantomsection\label{_Toc239502826}{}
all\_model\_metrics\_grouped\_bar}
\end{figure}


\section{Per-Class Performance of the Best
Model}\label{per-class-performance-of-the-best-model}

Performance is not uniform across categories for the proposed hybrid
model.
\begin{table}[htbp]
\centering
\caption{Per-Class Performance of the Proposed Hybrid Model}
\label{tab:per-class-performance}

\renewcommand{\arraystretch}{1.3}
\small

\begin{tabularx}{\textwidth}{
>{\raggedright\arraybackslash}X
>{\centering\arraybackslash}m{1.8cm}
>{\centering\arraybackslash}m{1.8cm}
>{\centering\arraybackslash}m{1.6cm}
>{\centering\arraybackslash}m{1.8cm}
}
\toprule

\textbf{Category} &
\textbf{Precision} &
\textbf{Recall} &
\textbf{F1} &
\textbf{Support} \\

\midrule

Data Breach / Leakage & 1.00 & 0.99 & 1.00 & 114 \\

DoS/DDoS Attack & 0.99 & 0.98 & 0.99 & 120 \\

Espionage / APT Campaign & 0.99 & 0.98 & 0.98 & 140 \\

Vulnerability Exploitation & 0.97 & 0.95 & 0.96 & 182 \\

Ransomware & 0.96 & 0.89 & 0.93 & 179 \\

Threat Actor Activity & 0.90 & 0.92 & 0.91 & 186 \\

Malware Infection & 0.89 & 0.97 & 0.92 & 263 \\

Phishing / Social Engineering & 0.89 & 0.92 & 0.90 & 119 \\

Other / Unspecified & 0.90 & 0.80 & 0.85 & 117 \\

\midrule

 
\textbf{Macro Average} &
\textbf{0.94} &
\textbf{0.93} &
\textbf{0.94} &
\textbf{1420} \\

 
\textbf{Weighted Average} &
\textbf{0.94} &
\textbf{0.94} &
\textbf{0.94} &
\textbf{1420} \\

\bottomrule
\end{tabularx}
\end{table}
The four categories with narrow, well-defined vocabulary, Data
Breach/Leakage, DoS/DDoS Attack, and
Vulnerability Exploitation, are classified almost perfectly (F1 between
0.96 and 1.00). The weakest category by a clear margin is
Other/Unspecified (F1 0.85, recall 0.80), which is expected given that
it is defined negatively, as whatever does not match any of the other
eight categories\textquotesingle{} keyword-and-lexicon rules, and
therefore carries no consistent internal vocabulary of its own for the
model to learn. 
\newline Malware Infection, the largest category by support,
shows a similar pattern: high recall but comparatively lower precision
(0.89), suggesting the model over-generates this label somewhat,
plausibly absorbing some records that a stricter rule would have placed
in a more specific category. The confusion is plausibly directional
rather than random: threat-actor and malware descriptions frequently
co-occur in the same narrative, a report naming both a group and the
tooling it deploys, which is exactly the kind of record the
lexicon-based label construction in Chapter 3 assigns to whichever
category\textquotesingle s pattern is checked first, and is a natural
source of residual ambiguity that a purely keyword-and-lexicon label
cannot fully resolve.

\begin{figure}
\centering
\includegraphics[width=6.5in,height=3.87361in]{images/media/image14.png}
\caption{\protect\phantomsection\label{_Toc239502827}{} per class
metrics visualized}
\end{figure}
\section{Diagnostic Error
Analysis}\label{diagnostic-error-analysis}

Out-of-fold predictions from the strongest locally-evaluated classical
baseline (the support vector machine) were cross-referenced against the
entity-aware measures from Section 3.5 to test whether misclassification
correlates with structural record properties. Accuracy on records with a
high Rare Entity Ratio was 0.815, compared with 0.903 for records with
no rare entities, an 8.8-percentage-point gap indicating that records
involving infrequently occurring entity types are measurably harder to
classify correctly. Accuracy segmented by Threat Context Score showed a
similar, non-monotonic pattern: 0.947 for low-context records, falling
to 0.828 for mid-context records and recovering slightly to 0.846 for
high-context records, suggesting that moderately ambiguous records,
rather than the most information-sparse ones, are where classical
baselines are weakest.
\section{Summary of Findings}

Three results anchor this chapter. The proposed hybrid model outperforms
every baseline evaluated, by a margin (2.73 macro-F1 points over the
best classical baseline) too large to attribute to cross-validation
noise given the reported standard deviations. The performance gain is
concentrated unevenly across categories, with near-perfect
classification of narrowly-defined categories and a clear, reproducible
weakness on the Other/Unspecified category specifically. That weakness
is independently corroborated by the entity-aware diagnostic analysis,
which identifies structurally similar records, those lacking a strong
lexical or structural signal, as the hardest case for the classical
baselines as well. Together, these three results support a single,
coherent account of the system\textquotesingle s behaviour, rather than
three unrelated observations.
\section{Summary}\label{summary-3}

In summary, Chapter 6 provides an exhaustive empirical evaluation of the
proposed SecureBERT-BiLSTM-Attention system. It confirms a statistically
significant performance gain over all reference baselines, validates the
domain-adaptation hypothesis, and identifies the exact boundary of the
system\textquotesingle s prediction boundaries through a structured
diagnostic error framework.

\chapter{Conclusion and Future Work}
\section{Conclusion}

This dissertation set out to build and evaluate a domain-adapted hybrid
Securebert-BiLSTM-Attention architecture for cyber threat intelligence
classification, motivated by a specific, literature-grounded gap: no
published hybrid transformer-recurrent CTI classifier had tested a
domain-adapted encoder in place of general-purpose BERT. Working toward
this, the project constructed and validated a genuine classification
label where none existed in the source data, identified and corrected a
label-construction dependency that would have undermined the
system\textquotesingle s real-world validity, curated and structured the
target dataset, and built a benchmarked pipeline spanning classical,
recurrent, and hybrid transformer architectures. The proposed hybrid
model achieved a macro-F1 of 0.9370, exceeding the strongest classical
baseline (TF-IDF with SVM, 0.9097) by 2.73 points and the published
literature benchmark for an equivalent general-BERT hybrid (Hidayatulloh
et
al.{[}\href{file:///D:/Downloads/final-report-formatted-v4(1).docx\#ref_10}{\ul{10}}{]})
by 7.7 points, directly confirming the domain-adaptation hypothesis this
dissertation set out to test.Beyond the model comparison itself, this
work contributes a diagnostic evaluation approach that links
entity-aware structural measures to model error patterns, converging
with the per-class results on the same underlying finding, that
structurally ambiguous, low-signal records remain the hardest case for
every model evaluated, and a lightweight, unsupervised
mitigation-recommendation step that extends the pipeline from
classification into a more directly actionable output.
\section{Discussion}

This project set out to test whether a hybrid
Securebert-BiLSTM-Attention architecture using a domain-adapted encoder
improves cyber threat intelligence classification beyond what classical
machine learning and a from-scratch BiLSTM achieve, and whether
entity-aware exploratory measures can meaningfully explain where such a
model struggles. Both questions are answered affirmatively by the
results in Chapter 6: the proposed hybrid model achieved a macro-F1 of
0.9370, exceeding the strongest classical baseline by 2.73 points and
the published general-BERT hybrid benchmark (Hidayatulloh et
al.{[}\href{file:///D:/Downloads/final-report-formatted-v4(1).docx\#ref_10}{\ul{10}}{]})
by 7.7 points, and the diagnostic analysis identified a specific,
consistent source of remaining difficulty, low-signal records without a
strong lexical or structural fingerprint, rather than leaving model
weaknesses undifferentiated behind a single aggregate score.Two findings
are particularly worth drawing out. First, dataset scale and class
balance measurably affected baseline performance independently of the
hybrid architecture itself: the BiLSTM\textquotesingle s performance was
sensitive to the training data size, showing that deep sequence models
trained from scratch on limited domain datasets face substantial
representation limitations, a constraint that the domain-adapted
pretraining of SecureBERT effectively bypasses. Second, the same
weakness surfaces from two independent directions: the Other/Unspecified
category, structurally the least lexically consistent of the nine, is
also the hybrid model\textquotesingle s weakest by per-class F1, and the
diagnostic analysis based on entity-aware measures shows the same
pattern of difficulty concentrated in structurally ambiguous records.
Two different analytical approaches converging on the same conclusion is
a stronger basis for that conclusion than either alone. Set against the
wider literature, the result also clarifies where the improvement
actually comes from. Hidayatulloh et
al.{[}\href{file:///D:/Downloads/final-report-formatted-v4(1).docx\#ref_10}{\ul{10}}{]}demonstrate
that a hybrid architecture improves on a standalone general-purpose
BERT; this project\textquotesingle s result, that substituting a
domain-adapted encoder into the same hybrid structure improves further
still, extends rather than merely replicates that finding, and directly
closes the specific gap identified in the literature review: no
published hybrid transformer-recurrent CTI classifier had previously
tested a domain-adapted encoder in place of general-purpose BERT.
\section{Limitations}

Several limitations should be weighed alongside these results. The
text-only label construction approach, while methodologically necessary
to avoid the label depending on information a deployed system would not
have, remains a keyword-and-lexicon heuristic rather than human
annotation, and the Other/Unspecified category, though substantially
reduced from its earlier size, is also the category where every model in
this comparison performs worst, a limitation of the label rather than of
any individual model. Finally, transformer fine-tuning was constrained
to a small number of epochs and a frozen lower-encoder configuration,
chosen specifically to manage overfitting risk at this dataset size
rather than to maximise raw performance; whether a larger real dataset
would support a less constrained fine-tuning regime, and whether that
would widen or narrow the margin over the classical baselines, remains
an open question this project\textquotesingle s scope does not resolve.
\section{Practical CTI
Triage Deployment Pattern}

For practical CTI triage, these results suggest a specific deployment
pattern rather than a blanket recommendation: the proposed model is
highly reliable on well-defined threat categories, sufficiently so that
its predictions on Data Breach/Leakage, DoS/DDoS Attack, Espionage/APT
Campaign, and Vulnerability Exploitation could plausibly support
automated or lightly-reviewed triage, while its comparatively weaker
performance on ambiguous, low-signal records suggests that predictions
falling into or near the Other/Unspecified category, or records
independently flagged by the diagnostic measures as high in Rare Entity
Ratio or moderate in Threat Context Score, are better routed to human
review rather than automated action. This is a more actionable
operational recommendation than a single aggregate accuracy figure would
support, and follows directly from treating per-class and diagnostic
results as first-class findings rather than as a secondary breakdown of
a headline number.
\section{Reflection on
Project Development}

The development of this project highlighted the critical importance of a
structured software and methodology workflow. Iterating on the priority
rules for label construction and discovering text formatting
inconsistencies earlier prevented cascading errors in model fine-tuning.
The rigorous logging and saving of intermediate datasets proved to be an
invaluable engineering practice.
\section{Research
Contributions}

This research makes three key contributions:

1. Technical: Demonstrating that combining a domain-adapted encoder with
a sequence sequence model and attention compounds the individual
performance benefits of each approach on specialized threat narratives.

2. Methodological: Establishing a diagnostic evaluation framework that
uses pre-existing, low-cost entity annotations to systematically locate
and explain model classification performance.

3. Practical: Providing a practical and modular automated CTI
classification tool with explicit confidence bounds and clear escalation
criteria for human analyst review.
\section{Project
Struggles and Solutions}

The primary struggle during project development centered on handling
highly imbalanced threat category distributions within a limited sample
size. This was successfully resolved by applying stratified
cross-validation and balanced class weighting, which preserved
representation integrity and allowed the models to train without severe
overfitting or class-selection bias.
\section{Future Work}

Future work could extend this project in several directions: replacing
the keyword-and-lexicon label construction with a small amount of human
annotation to establish a genuine gold-standard subset against which the
constructed labels could be formally validated; exploring true
multi-label classification, since manual inspection during this project
confirmed that some records plausibly belong to more than one threat
category; acquiring additional real CTI records to expand the
representation of minority threat categories; and extending the
mitigation-recommendation step with a generative component that adapts a
retrieved solution to the specifics of the new record, rather than
returning it unmodified.
\section{Detailed Model
Comparison and Design Reasoning}

The comparative experimental design of this study was selected
specifically to isolate the exact source of performance improvements. By
starting with simple TF-IDF lexical models, progressing to from-scratch
BiLSTM architectures, and finally testing BERT and SecureBERT hybrids,
we verified that domain pretraining on specific cybersecurity corpora
contributes the largest incremental macro-F1 gain, which is then further
stabilized by the BiLSTM sequence layer.
\section{Feature
Importance Discussion}

Analyzing the attention weights of the hybrid proposed model provided a
direct secondary mechanism for verifying that the system was attending
to authentic threat-intelligence terminology rather than noise or stop
words. High attention scores consistently aligned with malware family
names, CVE identifiers, and actionable network vectors.
\section{Additional
Analytical Discussion}

An analytical examination of structural relations within the records
confirmed that threat narratives are highly relational in nature. Threat
actors are frequently linked to specific campaigns and malware vectors,
confirming that further integration of relational features is a highly
promising direction.
\section{Final
Validation Against Research Questions}

The results systematically answer the primary research question set out
in Section 1.7: substituting SecureBERT for general BERT within the
hybrid model does improve F1 performance by a large margin (+7.7
points), and the system significantly outperforms classical baselines
(+2.73 points). Supporting questions regarding target construction and
diagnostic annotations were similarly answered in the affirmative.
\section{Final
Technical Positioning}

This research positions itself at the intersection of deep sequence
modeling and domain-adapted pretraining. It offers a highly lightweight,
supervised alternative to ungrounded generative LLMs that balances high
predictive accuracy with direct auditability and operational
explainability.
\section{Summary}

Overall, Chapter 7 synthesizes the study\textquotesingle s findings,
highlighting that domain adaptation and hybrid neural architectures
represent a highly effective paradigm for automative cyber threat
classification, while setting out a clear agenda for future work.
\section{Final Word on
Technical Depth}

The rigorous and systematic verification of every stage in this pipeline
from text preprocessing to out-of-fold cross-validation and entity-aware
error diagnostics ensures that the findings of this dissertation are
highly robust, repeatable, and academically credible.
% \cleardoublepage
% \printbibliography[heading=bibintoc,title={Bibliography}]
\cleardoublepage
\FloatBarrier
\label{bib}
\nocite{*}
\printbibliography[
    heading=bibintoc,
    title={Bibliography}
]
% \section{References}\label{references}

% {[}1{]} G. Husari, E. Al-Shaer, M. Ahmed, B. Chu, and X. Niu, "TTPDrill:
% Automatic and Accurate Extraction of Threat Actions from Unstructured
% Text of CTI Sources," in Proceedings of the 33rd Annual Computer
% Security Applications Conference (ACSAC), 2017, pp. 103-115.
% \href{https://doi.org/10.1145/3134600.3134646}{\ul{doi:
% 10.1145/3134600.3134646}}

% \protect\phantomsection\label{ref_2}{}{[}2{]} M. Büchel et al., "SoK:
% Automated TTP Extraction from CTI Reports -- Are We There Yet?" in
% Proceedings of the 34th USENIX Security Symposium, 2025, pp. 4621-4641.
% {[}Online{]}. Available:
% \href{https://www.usenix.org/conference/usenixsecurity25/presentation/buechel}{\ul{https://www.usenix.org/conference/usenixsecurity25/presentation/buechel}}

% \protect\phantomsection\label{ref_3}{}{[}3{]} E. Aghaei, X. Niu, W.
% Shadid, and E. Al-Shaer, "SecureBERT: A Domain-Specific Language Model
% for Cybersecurity," arXiv preprint arXiv:2204.02685, 2022. {[}Online{]}.
% Available:
% \href{https://arxiv.org/abs/2204.02685}{\ul{https://arxiv.org/abs/2204.02685}}

% \protect\phantomsection\label{ref_4}{}{[}4{]} M. Bayer, P. Kuehn, R.
% Shanehsaz, and C. Reuter, "CySecBERT: A Domain-Adapted Language Model
% for the Cybersecurity Domain," arXiv preprint arXiv:2212.02974, 2022.
% {[}Online{]}. Available:
% \href{https://arxiv.org/abs/2212.02974}{\ul{https://arxiv.org/abs/2212.02974}}

% \protect\phantomsection\label{ref_5}{}{[}5{]} M. T. Alam, D. Bhusal, L.
% Nguyen, and N. Rastogi, "CTIBench: A Benchmark for Evaluating LLMs in
% Cyber Threat Intelligence," in Advances in Neural Information Processing
% Systems (NeurIPS 2024), Datasets and Benchmarks Track, 2024.
% {[}Online{]}. Available:
% \href{https://arxiv.org/abs/2406.07599}{\ul{https://arxiv.org/abs/2406.07599}}

% \protect\phantomsection\label{ref_6}{}{[}6{]} A. Vaswani et al.,
% "Attention Is All You Need," arXiv preprint arXiv:1706.03762, 2017.
% {[}Online{]}. Available:
% \href{https://arxiv.org/abs/1706.03762}{\ul{https://arxiv.org/abs/1706.03762}}

% \protect\phantomsection\label{ref_7}{}{[}7{]} J. Devlin, M. W. Chang, K.
% Lee, and K. Toutanova, "BERT: Pre-training of Deep Bidirectional
% Transformers for Language Understanding," arXiv preprint
% arXiv:1810.04805, 2018. {[}Online{]}. Available:
% \href{https://arxiv.org/abs/1810.04805}{\ul{https://arxiv.org/abs/1810.04805}}

% \protect\phantomsection\label{ref_8}{}{[}8{]} S. Hochreiter and J.
% Schmidhuber, "Long Short-Term Memory," Neural Computation, vol. 9, no.
% 8, pp. 1735-1780, 1997. \url{https://doi.org/10.1162/neco.1997.9.8.1735}

% \protect\phantomsection\label{ref_9}{}{[}9{]} S. Jiang, S. Zhao, K. Hou,
% Y. Liu, and L. Zhang, "A BERT-BiLSTM-CRF Model for Chinese Electronic
% Medical Records Named Entity Recognition," in Proceedings of the 12th
% International Conference on Intelligent Computation Technology and
% Automation (ICICTA), 2019, pp. 166-169.
% \href{https://doi.org/10.1109/ICICTA49267.2019.00043}{\ul{doi:
% 10.1109/ICICTA49267.2019.00043}}

% \protect\phantomsection\label{ref_10}{}{[}10{]} S. Hidayatulloh, S.
% Topiq, I. Hariyanti, and D. Sandini, "Hybrid BERT-BiLSTM Architecture
% for Enhanced Cyber Threat Intelligence Classification," International
% Journal of Advanced Computer Science and Applications, vol. 17, no. 2,
% pp. 305-315, 2026.
% \href{https://doi.org/10.14569/IJACSA.2026.0170233}{\ul{doi:
% 10.14569/IJACSA.2026.0170233}}

% \protect\phantomsection\label{ref_11}{}{[}11{]} G. Xiang, C. Shi, and Y.
% Zhang, "An APT Event Extraction Method Based on BERT-BiGRU-CRF for APT
% Attack Detection," Electronics, vol. 12, no. 15, p. 3349, 2023.
% \href{https://doi.org/10.3390/electronics12153349}{\ul{doi:
% 10.3390/electronics12153349}}

% \protect\phantomsection\label{ref_12}{}{[}12{]} Y. Liu and Y. Dai, "Deep
% Learning in Cybersecurity: A Hybrid BERT-LSTM Network for SQL Injection
% Attack Detection," IET Information Security, vol. 2024, p. 5565950,
% 2024. \href{https://doi.org/10.1049/2024/5565950}{\ul{doi:
% 10.1049/2024/5565950}}

% \protect\phantomsection\label{ref_13}{}{[}13{]} Z. Chen et al., ``Load
% Forecasting Based on LSTM Neural Network and Applicable to Loads of
% `Replacement of Coal with Electricity','' Journal of Electrical
% Engineering \& Technology, vol. 16, pp. 2333--2342, 2021.

% \url{https://doi.org/10.1007/s42835-021-00768-8}

% {[}14{]} M. A. Ferrag et al., "Revolutionizing Cyber Threat Detection
% with Large Language Models: A Privacy-Preserving BERT-Based Lightweight
% Model for IoT/IIoT Devices," IEEE Access, vol. 12, pp. 12345-12360,
% 2024. {[}Online{]}. Available:
% \href{https://arxiv.org/abs/2306.14263}{\ul{https://arxiv.org/abs/2306.14263}}

% {[}15{]} P. Satyapanich, F. Ferraro, and T. Finin, "CASIE: Extracting
% Cybersecurity Event Information from Text," in Proceedings of the AAAI
% Conference on Artificial Intelligence, vol. 34, no. 05, pp. 8749-8756,
% 2020. \href{https://doi.org/10.1609/aaai.v34i05.6401}{\ul{doi:
% 10.1609/aaai.v34i05.6401}}

% {[}16{]} P. Ranade, A. Piplai, A. Joshi, and T. Finin, "CyBERT:
% Contextualized Embeddings for the Cybersecurity Domain," in Proceedings
% of the 2021 IEEE International Conference on Big Data (Big Data), 2021,
% pp. 3334-3341.

% \url{https://doi.org/10.1109/BigData52589.2021.9671824}

% {[}17{]} E. Aghaei et al., "SecureBERT 2.0: Advanced Language Model for
% Cybersecurity Intelligence," arXiv preprint arXiv:2510.00240, 2025.
% {[}Online{]}. Available: \url{https://arxiv.org/abs/2510.00240}

% {[}18{]} J. Kaur, M. Prabha, M. Samiun, S. N. Hasan, and R. Hasan,
% "Comparative Analysis of Transformer and LSTM Architectures for
% Cybersecurity Threat Detection Using Machine Learning," EAI Endorsed
% Transactions on AI and Robotics, vol. 8, p. 9759, 2024.\\
% \url{https://doi.org/10.4108/airo.9759}

% {[}19{]} J. Park and H. You, "Cyber-Attack Technique Classification
% Using Two-Stage Trained Large Language Models," arXiv preprint
% arXiv:2411.18755, 2024. {[}Online{]}. Available:
% \href{https://arxiv.org/abs/2411.18755}{\ul{https://arxiv.org/abs/2411.18755}}

% {[}20{]} L. Arndt et al., "ThreatCrawl: A BERT-based Focused Crawler for
% the Cybersecurity Domain," arXiv preprint arXiv:2304.11960, 2023.
% {[}Online{]}. Available:
% \href{https://arxiv.org/abs/2304.11960}{\ul{https://arxiv.org/abs/2304.11960}}
\appendix
\chapter{Appendix}

\section{Streamlit Dashboard
}\label{a.1-streamlit-dashboard-}

The custom electricity load forecasting dashboard has been successfully
compiled and deployed using Streamlit.

\vspace{0.5cm}

\begin{figure}[H]
    \centering
    \includegraphics[width=\textwidth]{images/media/image15.png}
    \caption{User Login and
Authentication Panel CYBER threat classifier Streamlit App}
    \label{_Toc239462872}
\end{figure}
Above Figure presents the login/signup interface of the developed Streamlit-based Cyber
Threat Intelligence (CTI) dashboard. The interface allows users authenticate to the
dashboard and then enter cybersecurity-related text into a dedicated text area before
initiating the classification process.

\begin{figure}[H]
    \centering
    \includegraphics[width=\textwidth]{images/media/image16.png}
    \caption{Enter input CTI text in streamlit dashboard}
    \label{_Toc239462873}
\end{figure}
Above Figure presents the input interface of the developed Streamlit-based Cyber Threat
Intelligence (CTI) dashboard. The interface allows authenticated users to enter cybersecurity-
related text into a dedicated text area before initiating the classification process. Once
the Classify Threat button is selected, the input text is tokenised using the SecureBERT
tokenizer and prepared for inference by the deployed Hybrid SecureBERT–BiLSTM–
Attention model. The dashboard was designed to provide a simple and interactive in-
terface for demonstrating the practical deployment of the proposed threat classification
system.
\clearpage

\begin{figure}[H]
    \centering
    \includegraphics[width=\textwidth]{images/media/image17.png}
    \caption{Prediction result of classified threat}
    \label{_Toc239462874}
\end{figure}

Above Figure illustrates the prediction output generated by the deployed Streamlit dash-
board after processing a CTI record. The dashboard displays the predicted threat cate-
gory, a confidence score, and a visual confidence indicator to improve result interpretation.
During inference, the input text passes through the SecureBERT encoder to generate con-
textual representations, which are subsequently processed by the BiLSTM and Attention
layers before the final classification layer produces the predicted threat category. en-
abling users to classify CTI text through an interactive web application rather than a
notebook-based environment
\clearpage

\end{document}
% --------------